\documentclass[11pt]{article}

    \usepackage[breakable]{tcolorbox}
    \usepackage{parskip} % Stop auto-indenting (to mimic markdown behaviour)
    

    % Basic figure setup, for now with no caption control since it's done
    % automatically by Pandoc (which extracts ![](path) syntax from Markdown).
    \usepackage{graphicx}
    % Maintain compatibility with old templates. Remove in nbconvert 6.0
    \let\Oldincludegraphics\includegraphics
    % Ensure that by default, figures have no caption (until we provide a
    % proper Figure object with a Caption API and a way to capture that
    % in the conversion process - todo).
    \usepackage{caption}
    \DeclareCaptionFormat{nocaption}{}
    \captionsetup{format=nocaption,aboveskip=0pt,belowskip=0pt}

    \usepackage{float}
    \floatplacement{figure}{H} % forces figures to be placed at the correct location
    \usepackage{xcolor} % Allow colors to be defined
    \usepackage{enumerate} % Needed for markdown enumerations to work
    \usepackage{geometry} % Used to adjust the document margins
    \usepackage{amsmath} % Equations
    \usepackage{amssymb} % Equations
    \usepackage{textcomp} % defines textquotesingle
    % Hack from http://tex.stackexchange.com/a/47451/13684:
    \AtBeginDocument{%
        \def\PYZsq{\textquotesingle}% Upright quotes in Pygmentized code
    }
    \usepackage{upquote} % Upright quotes for verbatim code
    \usepackage{eurosym} % defines \euro

    \usepackage{iftex}
    \ifPDFTeX
        \usepackage[T1]{fontenc}
        \IfFileExists{alphabeta.sty}{
              \usepackage{alphabeta}
          }{
              \usepackage[mathletters]{ucs}
              \usepackage[utf8x]{inputenc}
          }
    \else
        \usepackage{fontspec}
        \usepackage{unicode-math}
    \fi

    \usepackage{fancyvrb} % verbatim replacement that allows latex
    \usepackage{grffile} % extends the file name processing of package graphics
                         % to support a larger range
    \makeatletter % fix for old versions of grffile with XeLaTeX
    \@ifpackagelater{grffile}{2019/11/01}
    {
      % Do nothing on new versions
    }
    {
      \def\Gread@@xetex#1{%
        \IfFileExists{"\Gin@base".bb}%
        {\Gread@eps{\Gin@base.bb}}%
        {\Gread@@xetex@aux#1}%
      }
    }
    \makeatother
    \usepackage[Export]{adjustbox} % Used to constrain images to a maximum size
    \adjustboxset{max size={0.9\linewidth}{0.9\paperheight}}

    % The hyperref package gives us a pdf with properly built
    % internal navigation ('pdf bookmarks' for the table of contents,
    % internal cross-reference links, web links for URLs, etc.)
    \usepackage{hyperref}
    % The default LaTeX title has an obnoxious amount of whitespace. By default,
    % titling removes some of it. It also provides customization options.
    \usepackage{titling}
    \usepackage{longtable} % longtable support required by pandoc >1.10
    \usepackage{booktabs}  % table support for pandoc > 1.12.2
    \usepackage{array}     % table support for pandoc >= 2.11.3
    \usepackage{calc}      % table minipage width calculation for pandoc >= 2.11.1
    \usepackage[inline]{enumitem} % IRkernel/repr support (it uses the enumerate* environment)
    \usepackage[normalem]{ulem} % ulem is needed to support strikethroughs (\sout)
                                % normalem makes italics be italics, not underlines
    \usepackage{mathrsfs}
    

    
    % Colors for the hyperref package
    \definecolor{urlcolor}{rgb}{0,.145,.698}
    \definecolor{linkcolor}{rgb}{.71,0.21,0.01}
    \definecolor{citecolor}{rgb}{.12,.54,.11}

    % ANSI colors
    \definecolor{ansi-black}{HTML}{3E424D}
    \definecolor{ansi-black-intense}{HTML}{282C36}
    \definecolor{ansi-red}{HTML}{E75C58}
    \definecolor{ansi-red-intense}{HTML}{B22B31}
    \definecolor{ansi-green}{HTML}{00A250}
    \definecolor{ansi-green-intense}{HTML}{007427}
    \definecolor{ansi-yellow}{HTML}{DDB62B}
    \definecolor{ansi-yellow-intense}{HTML}{B27D12}
    \definecolor{ansi-blue}{HTML}{208FFB}
    \definecolor{ansi-blue-intense}{HTML}{0065CA}
    \definecolor{ansi-magenta}{HTML}{D160C4}
    \definecolor{ansi-magenta-intense}{HTML}{A03196}
    \definecolor{ansi-cyan}{HTML}{60C6C8}
    \definecolor{ansi-cyan-intense}{HTML}{258F8F}
    \definecolor{ansi-white}{HTML}{C5C1B4}
    \definecolor{ansi-white-intense}{HTML}{A1A6B2}
    \definecolor{ansi-default-inverse-fg}{HTML}{FFFFFF}
    \definecolor{ansi-default-inverse-bg}{HTML}{000000}

    % common color for the border for error outputs.
    \definecolor{outerrorbackground}{HTML}{FFDFDF}

    % commands and environments needed by pandoc snippets
    % extracted from the output of `pandoc -s`
    \providecommand{\tightlist}{%
      \setlength{\itemsep}{0pt}\setlength{\parskip}{0pt}}
    \DefineVerbatimEnvironment{Highlighting}{Verbatim}{commandchars=\\\{\}}
    % Add ',fontsize=\small' for more characters per line
    \newenvironment{Shaded}{}{}
    \newcommand{\KeywordTok}[1]{\textcolor[rgb]{0.00,0.44,0.13}{\textbf{{#1}}}}
    \newcommand{\DataTypeTok}[1]{\textcolor[rgb]{0.56,0.13,0.00}{{#1}}}
    \newcommand{\DecValTok}[1]{\textcolor[rgb]{0.25,0.63,0.44}{{#1}}}
    \newcommand{\BaseNTok}[1]{\textcolor[rgb]{0.25,0.63,0.44}{{#1}}}
    \newcommand{\FloatTok}[1]{\textcolor[rgb]{0.25,0.63,0.44}{{#1}}}
    \newcommand{\CharTok}[1]{\textcolor[rgb]{0.25,0.44,0.63}{{#1}}}
    \newcommand{\StringTok}[1]{\textcolor[rgb]{0.25,0.44,0.63}{{#1}}}
    \newcommand{\CommentTok}[1]{\textcolor[rgb]{0.38,0.63,0.69}{\textit{{#1}}}}
    \newcommand{\OtherTok}[1]{\textcolor[rgb]{0.00,0.44,0.13}{{#1}}}
    \newcommand{\AlertTok}[1]{\textcolor[rgb]{1.00,0.00,0.00}{\textbf{{#1}}}}
    \newcommand{\FunctionTok}[1]{\textcolor[rgb]{0.02,0.16,0.49}{{#1}}}
    \newcommand{\RegionMarkerTok}[1]{{#1}}
    \newcommand{\ErrorTok}[1]{\textcolor[rgb]{1.00,0.00,0.00}{\textbf{{#1}}}}
    \newcommand{\NormalTok}[1]{{#1}}

    % Additional commands for more recent versions of Pandoc
    \newcommand{\ConstantTok}[1]{\textcolor[rgb]{0.53,0.00,0.00}{{#1}}}
    \newcommand{\SpecialCharTok}[1]{\textcolor[rgb]{0.25,0.44,0.63}{{#1}}}
    \newcommand{\VerbatimStringTok}[1]{\textcolor[rgb]{0.25,0.44,0.63}{{#1}}}
    \newcommand{\SpecialStringTok}[1]{\textcolor[rgb]{0.73,0.40,0.53}{{#1}}}
    \newcommand{\ImportTok}[1]{{#1}}
    \newcommand{\DocumentationTok}[1]{\textcolor[rgb]{0.73,0.13,0.13}{\textit{{#1}}}}
    \newcommand{\AnnotationTok}[1]{\textcolor[rgb]{0.38,0.63,0.69}{\textbf{\textit{{#1}}}}}
    \newcommand{\CommentVarTok}[1]{\textcolor[rgb]{0.38,0.63,0.69}{\textbf{\textit{{#1}}}}}
    \newcommand{\VariableTok}[1]{\textcolor[rgb]{0.10,0.09,0.49}{{#1}}}
    \newcommand{\ControlFlowTok}[1]{\textcolor[rgb]{0.00,0.44,0.13}{\textbf{{#1}}}}
    \newcommand{\OperatorTok}[1]{\textcolor[rgb]{0.40,0.40,0.40}{{#1}}}
    \newcommand{\BuiltInTok}[1]{{#1}}
    \newcommand{\ExtensionTok}[1]{{#1}}
    \newcommand{\PreprocessorTok}[1]{\textcolor[rgb]{0.74,0.48,0.00}{{#1}}}
    \newcommand{\AttributeTok}[1]{\textcolor[rgb]{0.49,0.56,0.16}{{#1}}}
    \newcommand{\InformationTok}[1]{\textcolor[rgb]{0.38,0.63,0.69}{\textbf{\textit{{#1}}}}}
    \newcommand{\WarningTok}[1]{\textcolor[rgb]{0.38,0.63,0.69}{\textbf{\textit{{#1}}}}}


    % Define a nice break command that doesn't care if a line doesn't already
    % exist.
    \def\br{\hspace*{\fill} \\* }
    % Math Jax compatibility definitions
    \def\gt{>}
    \def\lt{<}
    \let\Oldtex\TeX
    \let\Oldlatex\LaTeX
    \renewcommand{\TeX}{\textrm{\Oldtex}}
    \renewcommand{\LaTeX}{\textrm{\Oldlatex}}
    % Document parameters
    % Document title
    \title{hw-02}
    
    
    
    
    
% Pygments definitions
\makeatletter
\def\PY@reset{\let\PY@it=\relax \let\PY@bf=\relax%
    \let\PY@ul=\relax \let\PY@tc=\relax%
    \let\PY@bc=\relax \let\PY@ff=\relax}
\def\PY@tok#1{\csname PY@tok@#1\endcsname}
\def\PY@toks#1+{\ifx\relax#1\empty\else%
    \PY@tok{#1}\expandafter\PY@toks\fi}
\def\PY@do#1{\PY@bc{\PY@tc{\PY@ul{%
    \PY@it{\PY@bf{\PY@ff{#1}}}}}}}
\def\PY#1#2{\PY@reset\PY@toks#1+\relax+\PY@do{#2}}

\@namedef{PY@tok@w}{\def\PY@tc##1{\textcolor[rgb]{0.73,0.73,0.73}{##1}}}
\@namedef{PY@tok@c}{\let\PY@it=\textit\def\PY@tc##1{\textcolor[rgb]{0.24,0.48,0.48}{##1}}}
\@namedef{PY@tok@cp}{\def\PY@tc##1{\textcolor[rgb]{0.61,0.40,0.00}{##1}}}
\@namedef{PY@tok@k}{\let\PY@bf=\textbf\def\PY@tc##1{\textcolor[rgb]{0.00,0.50,0.00}{##1}}}
\@namedef{PY@tok@kp}{\def\PY@tc##1{\textcolor[rgb]{0.00,0.50,0.00}{##1}}}
\@namedef{PY@tok@kt}{\def\PY@tc##1{\textcolor[rgb]{0.69,0.00,0.25}{##1}}}
\@namedef{PY@tok@o}{\def\PY@tc##1{\textcolor[rgb]{0.40,0.40,0.40}{##1}}}
\@namedef{PY@tok@ow}{\let\PY@bf=\textbf\def\PY@tc##1{\textcolor[rgb]{0.67,0.13,1.00}{##1}}}
\@namedef{PY@tok@nb}{\def\PY@tc##1{\textcolor[rgb]{0.00,0.50,0.00}{##1}}}
\@namedef{PY@tok@nf}{\def\PY@tc##1{\textcolor[rgb]{0.00,0.00,1.00}{##1}}}
\@namedef{PY@tok@nc}{\let\PY@bf=\textbf\def\PY@tc##1{\textcolor[rgb]{0.00,0.00,1.00}{##1}}}
\@namedef{PY@tok@nn}{\let\PY@bf=\textbf\def\PY@tc##1{\textcolor[rgb]{0.00,0.00,1.00}{##1}}}
\@namedef{PY@tok@ne}{\let\PY@bf=\textbf\def\PY@tc##1{\textcolor[rgb]{0.80,0.25,0.22}{##1}}}
\@namedef{PY@tok@nv}{\def\PY@tc##1{\textcolor[rgb]{0.10,0.09,0.49}{##1}}}
\@namedef{PY@tok@no}{\def\PY@tc##1{\textcolor[rgb]{0.53,0.00,0.00}{##1}}}
\@namedef{PY@tok@nl}{\def\PY@tc##1{\textcolor[rgb]{0.46,0.46,0.00}{##1}}}
\@namedef{PY@tok@ni}{\let\PY@bf=\textbf\def\PY@tc##1{\textcolor[rgb]{0.44,0.44,0.44}{##1}}}
\@namedef{PY@tok@na}{\def\PY@tc##1{\textcolor[rgb]{0.41,0.47,0.13}{##1}}}
\@namedef{PY@tok@nt}{\let\PY@bf=\textbf\def\PY@tc##1{\textcolor[rgb]{0.00,0.50,0.00}{##1}}}
\@namedef{PY@tok@nd}{\def\PY@tc##1{\textcolor[rgb]{0.67,0.13,1.00}{##1}}}
\@namedef{PY@tok@s}{\def\PY@tc##1{\textcolor[rgb]{0.73,0.13,0.13}{##1}}}
\@namedef{PY@tok@sd}{\let\PY@it=\textit\def\PY@tc##1{\textcolor[rgb]{0.73,0.13,0.13}{##1}}}
\@namedef{PY@tok@si}{\let\PY@bf=\textbf\def\PY@tc##1{\textcolor[rgb]{0.64,0.35,0.47}{##1}}}
\@namedef{PY@tok@se}{\let\PY@bf=\textbf\def\PY@tc##1{\textcolor[rgb]{0.67,0.36,0.12}{##1}}}
\@namedef{PY@tok@sr}{\def\PY@tc##1{\textcolor[rgb]{0.64,0.35,0.47}{##1}}}
\@namedef{PY@tok@ss}{\def\PY@tc##1{\textcolor[rgb]{0.10,0.09,0.49}{##1}}}
\@namedef{PY@tok@sx}{\def\PY@tc##1{\textcolor[rgb]{0.00,0.50,0.00}{##1}}}
\@namedef{PY@tok@m}{\def\PY@tc##1{\textcolor[rgb]{0.40,0.40,0.40}{##1}}}
\@namedef{PY@tok@gh}{\let\PY@bf=\textbf\def\PY@tc##1{\textcolor[rgb]{0.00,0.00,0.50}{##1}}}
\@namedef{PY@tok@gu}{\let\PY@bf=\textbf\def\PY@tc##1{\textcolor[rgb]{0.50,0.00,0.50}{##1}}}
\@namedef{PY@tok@gd}{\def\PY@tc##1{\textcolor[rgb]{0.63,0.00,0.00}{##1}}}
\@namedef{PY@tok@gi}{\def\PY@tc##1{\textcolor[rgb]{0.00,0.52,0.00}{##1}}}
\@namedef{PY@tok@gr}{\def\PY@tc##1{\textcolor[rgb]{0.89,0.00,0.00}{##1}}}
\@namedef{PY@tok@ge}{\let\PY@it=\textit}
\@namedef{PY@tok@gs}{\let\PY@bf=\textbf}
\@namedef{PY@tok@gp}{\let\PY@bf=\textbf\def\PY@tc##1{\textcolor[rgb]{0.00,0.00,0.50}{##1}}}
\@namedef{PY@tok@go}{\def\PY@tc##1{\textcolor[rgb]{0.44,0.44,0.44}{##1}}}
\@namedef{PY@tok@gt}{\def\PY@tc##1{\textcolor[rgb]{0.00,0.27,0.87}{##1}}}
\@namedef{PY@tok@err}{\def\PY@bc##1{{\setlength{\fboxsep}{\string -\fboxrule}\fcolorbox[rgb]{1.00,0.00,0.00}{1,1,1}{\strut ##1}}}}
\@namedef{PY@tok@kc}{\let\PY@bf=\textbf\def\PY@tc##1{\textcolor[rgb]{0.00,0.50,0.00}{##1}}}
\@namedef{PY@tok@kd}{\let\PY@bf=\textbf\def\PY@tc##1{\textcolor[rgb]{0.00,0.50,0.00}{##1}}}
\@namedef{PY@tok@kn}{\let\PY@bf=\textbf\def\PY@tc##1{\textcolor[rgb]{0.00,0.50,0.00}{##1}}}
\@namedef{PY@tok@kr}{\let\PY@bf=\textbf\def\PY@tc##1{\textcolor[rgb]{0.00,0.50,0.00}{##1}}}
\@namedef{PY@tok@bp}{\def\PY@tc##1{\textcolor[rgb]{0.00,0.50,0.00}{##1}}}
\@namedef{PY@tok@fm}{\def\PY@tc##1{\textcolor[rgb]{0.00,0.00,1.00}{##1}}}
\@namedef{PY@tok@vc}{\def\PY@tc##1{\textcolor[rgb]{0.10,0.09,0.49}{##1}}}
\@namedef{PY@tok@vg}{\def\PY@tc##1{\textcolor[rgb]{0.10,0.09,0.49}{##1}}}
\@namedef{PY@tok@vi}{\def\PY@tc##1{\textcolor[rgb]{0.10,0.09,0.49}{##1}}}
\@namedef{PY@tok@vm}{\def\PY@tc##1{\textcolor[rgb]{0.10,0.09,0.49}{##1}}}
\@namedef{PY@tok@sa}{\def\PY@tc##1{\textcolor[rgb]{0.73,0.13,0.13}{##1}}}
\@namedef{PY@tok@sb}{\def\PY@tc##1{\textcolor[rgb]{0.73,0.13,0.13}{##1}}}
\@namedef{PY@tok@sc}{\def\PY@tc##1{\textcolor[rgb]{0.73,0.13,0.13}{##1}}}
\@namedef{PY@tok@dl}{\def\PY@tc##1{\textcolor[rgb]{0.73,0.13,0.13}{##1}}}
\@namedef{PY@tok@s2}{\def\PY@tc##1{\textcolor[rgb]{0.73,0.13,0.13}{##1}}}
\@namedef{PY@tok@sh}{\def\PY@tc##1{\textcolor[rgb]{0.73,0.13,0.13}{##1}}}
\@namedef{PY@tok@s1}{\def\PY@tc##1{\textcolor[rgb]{0.73,0.13,0.13}{##1}}}
\@namedef{PY@tok@mb}{\def\PY@tc##1{\textcolor[rgb]{0.40,0.40,0.40}{##1}}}
\@namedef{PY@tok@mf}{\def\PY@tc##1{\textcolor[rgb]{0.40,0.40,0.40}{##1}}}
\@namedef{PY@tok@mh}{\def\PY@tc##1{\textcolor[rgb]{0.40,0.40,0.40}{##1}}}
\@namedef{PY@tok@mi}{\def\PY@tc##1{\textcolor[rgb]{0.40,0.40,0.40}{##1}}}
\@namedef{PY@tok@il}{\def\PY@tc##1{\textcolor[rgb]{0.40,0.40,0.40}{##1}}}
\@namedef{PY@tok@mo}{\def\PY@tc##1{\textcolor[rgb]{0.40,0.40,0.40}{##1}}}
\@namedef{PY@tok@ch}{\let\PY@it=\textit\def\PY@tc##1{\textcolor[rgb]{0.24,0.48,0.48}{##1}}}
\@namedef{PY@tok@cm}{\let\PY@it=\textit\def\PY@tc##1{\textcolor[rgb]{0.24,0.48,0.48}{##1}}}
\@namedef{PY@tok@cpf}{\let\PY@it=\textit\def\PY@tc##1{\textcolor[rgb]{0.24,0.48,0.48}{##1}}}
\@namedef{PY@tok@c1}{\let\PY@it=\textit\def\PY@tc##1{\textcolor[rgb]{0.24,0.48,0.48}{##1}}}
\@namedef{PY@tok@cs}{\let\PY@it=\textit\def\PY@tc##1{\textcolor[rgb]{0.24,0.48,0.48}{##1}}}

\def\PYZbs{\char`\\}
\def\PYZus{\char`\_}
\def\PYZob{\char`\{}
\def\PYZcb{\char`\}}
\def\PYZca{\char`\^}
\def\PYZam{\char`\&}
\def\PYZlt{\char`\<}
\def\PYZgt{\char`\>}
\def\PYZsh{\char`\#}
\def\PYZpc{\char`\%}
\def\PYZdl{\char`\$}
\def\PYZhy{\char`\-}
\def\PYZsq{\char`\'}
\def\PYZdq{\char`\"}
\def\PYZti{\char`\~}
% for compatibility with earlier versions
\def\PYZat{@}
\def\PYZlb{[}
\def\PYZrb{]}
\makeatother


    % For linebreaks inside Verbatim environment from package fancyvrb.
    \makeatletter
        \newbox\Wrappedcontinuationbox
        \newbox\Wrappedvisiblespacebox
        \newcommand*\Wrappedvisiblespace {\textcolor{red}{\textvisiblespace}}
        \newcommand*\Wrappedcontinuationsymbol {\textcolor{red}{\llap{\tiny$\m@th\hookrightarrow$}}}
        \newcommand*\Wrappedcontinuationindent {3ex }
        \newcommand*\Wrappedafterbreak {\kern\Wrappedcontinuationindent\copy\Wrappedcontinuationbox}
        % Take advantage of the already applied Pygments mark-up to insert
        % potential linebreaks for TeX processing.
        %        {, <, #, %, $, ' and ": go to next line.
        %        _, }, ^, &, >, - and ~: stay at end of broken line.
        % Use of \textquotesingle for straight quote.
        \newcommand*\Wrappedbreaksatspecials {%
            \def\PYGZus{\discretionary{\char`\_}{\Wrappedafterbreak}{\char`\_}}%
            \def\PYGZob{\discretionary{}{\Wrappedafterbreak\char`\{}{\char`\{}}%
            \def\PYGZcb{\discretionary{\char`\}}{\Wrappedafterbreak}{\char`\}}}%
            \def\PYGZca{\discretionary{\char`\^}{\Wrappedafterbreak}{\char`\^}}%
            \def\PYGZam{\discretionary{\char`\&}{\Wrappedafterbreak}{\char`\&}}%
            \def\PYGZlt{\discretionary{}{\Wrappedafterbreak\char`\<}{\char`\<}}%
            \def\PYGZgt{\discretionary{\char`\>}{\Wrappedafterbreak}{\char`\>}}%
            \def\PYGZsh{\discretionary{}{\Wrappedafterbreak\char`\#}{\char`\#}}%
            \def\PYGZpc{\discretionary{}{\Wrappedafterbreak\char`\%}{\char`\%}}%
            \def\PYGZdl{\discretionary{}{\Wrappedafterbreak\char`\$}{\char`\$}}%
            \def\PYGZhy{\discretionary{\char`\-}{\Wrappedafterbreak}{\char`\-}}%
            \def\PYGZsq{\discretionary{}{\Wrappedafterbreak\textquotesingle}{\textquotesingle}}%
            \def\PYGZdq{\discretionary{}{\Wrappedafterbreak\char`\"}{\char`\"}}%
            \def\PYGZti{\discretionary{\char`\~}{\Wrappedafterbreak}{\char`\~}}%
        }
        % Some characters . , ; ? ! / are not pygmentized.
        % This macro makes them "active" and they will insert potential linebreaks
        \newcommand*\Wrappedbreaksatpunct {%
            \lccode`\~`\.\lowercase{\def~}{\discretionary{\hbox{\char`\.}}{\Wrappedafterbreak}{\hbox{\char`\.}}}%
            \lccode`\~`\,\lowercase{\def~}{\discretionary{\hbox{\char`\,}}{\Wrappedafterbreak}{\hbox{\char`\,}}}%
            \lccode`\~`\;\lowercase{\def~}{\discretionary{\hbox{\char`\;}}{\Wrappedafterbreak}{\hbox{\char`\;}}}%
            \lccode`\~`\:\lowercase{\def~}{\discretionary{\hbox{\char`\:}}{\Wrappedafterbreak}{\hbox{\char`\:}}}%
            \lccode`\~`\?\lowercase{\def~}{\discretionary{\hbox{\char`\?}}{\Wrappedafterbreak}{\hbox{\char`\?}}}%
            \lccode`\~`\!\lowercase{\def~}{\discretionary{\hbox{\char`\!}}{\Wrappedafterbreak}{\hbox{\char`\!}}}%
            \lccode`\~`\/\lowercase{\def~}{\discretionary{\hbox{\char`\/}}{\Wrappedafterbreak}{\hbox{\char`\/}}}%
            \catcode`\.\active
            \catcode`\,\active
            \catcode`\;\active
            \catcode`\:\active
            \catcode`\?\active
            \catcode`\!\active
            \catcode`\/\active
            \lccode`\~`\~
        }
    \makeatother

    \let\OriginalVerbatim=\Verbatim
    \makeatletter
    \renewcommand{\Verbatim}[1][1]{%
        %\parskip\z@skip
        \sbox\Wrappedcontinuationbox {\Wrappedcontinuationsymbol}%
        \sbox\Wrappedvisiblespacebox {\FV@SetupFont\Wrappedvisiblespace}%
        \def\FancyVerbFormatLine ##1{\hsize\linewidth
            \vtop{\raggedright\hyphenpenalty\z@\exhyphenpenalty\z@
                \doublehyphendemerits\z@\finalhyphendemerits\z@
                \strut ##1\strut}%
        }%
        % If the linebreak is at a space, the latter will be displayed as visible
        % space at end of first line, and a continuation symbol starts next line.
        % Stretch/shrink are however usually zero for typewriter font.
        \def\FV@Space {%
            \nobreak\hskip\z@ plus\fontdimen3\font minus\fontdimen4\font
            \discretionary{\copy\Wrappedvisiblespacebox}{\Wrappedafterbreak}
            {\kern\fontdimen2\font}%
        }%

        % Allow breaks at special characters using \PYG... macros.
        \Wrappedbreaksatspecials
        % Breaks at punctuation characters . , ; ? ! and / need catcode=\active
        \OriginalVerbatim[#1,codes*=\Wrappedbreaksatpunct]%
    }
    \makeatother

    % Exact colors from NB
    \definecolor{incolor}{HTML}{303F9F}
    \definecolor{outcolor}{HTML}{D84315}
    \definecolor{cellborder}{HTML}{CFCFCF}
    \definecolor{cellbackground}{HTML}{F7F7F7}

    % prompt
    \makeatletter
    \newcommand{\boxspacing}{\kern\kvtcb@left@rule\kern\kvtcb@boxsep}
    \makeatother
    \newcommand{\prompt}[4]{
        {\ttfamily\llap{{\color{#2}[#3]:\hspace{3pt}#4}}\vspace{-\baselineskip}}
    }
    

    
    % Prevent overflowing lines due to hard-to-break entities
    \sloppy
    % Setup hyperref package
    \hypersetup{
      breaklinks=true,  % so long urls are correctly broken across lines
      colorlinks=true,
      urlcolor=urlcolor,
      linkcolor=linkcolor,
      citecolor=citecolor,
      }
    % Slightly bigger margins than the latex defaults
    
    \geometry{verbose,tmargin=1in,bmargin=1in,lmargin=1in,rmargin=1in}
    
    

\begin{document}
    
    \maketitle
    
    

    
    \hypertarget{homework-assignment-2}{%
\section{Homework assignment \#2}\label{homework-assignment-2}}

\textbf{YOUR NAME}

You are free to use any code developed in the class or the textbook.

When your solution contains computer code, please submit the work as a
Jupyter notebook. Please write your name near the top of the notebook.

If your solution contains hand-written calculations, it is probably
easiest to submit these on paper in class. You can type mathematics
using LaTeX expressions between dollar signs (if you know how to do
that), but hand written is usually easier and is fine.

All question numbers refer to exercises in the textbook.

    \hypertarget{section-2.1-1}{%
\subsection{Section 2.1 \#1}\label{section-2.1-1}}

Hint: here is some code to create a matrix based on a function to be
evaluated for each entry (i,j).

\textbf{Answers:}

\begin{enumerate}
\def\labelenumi{(\alph{enumi})}
\item
  \(5\) points are given, hence the maximum necessary degree of this
  polynomial is \(4\).
\item
  Interpolating polynomial
  \(p(t)=c_0+c_1 t + c_2 t^2 + c_3 t^3 + c_4 t^4\), corresponding linear
  system:
\end{enumerate}

\[\left\{\begin{array}{l}
c_0-c_1+c_2-c_3+c_4=0\\
c_0=1\\
c_0+2c_1+4c_2+8c_3+16c_4=0\\
c_0+3c_1+9c_2+27c_3+81c_4=1\\
c_0+4c_1+16c_2+64c_3+256c_4=2\\
\end{array}\right.\qquad
\Leftrightarrow\qquad
\begin{bmatrix}
1&-1&1&-1&1\\
1&0&0&0&0\\
1&2&4&8&16\\
1&3&9&27&81\\
1&4&16&64&256
\end{bmatrix}
\begin{bmatrix}c_0\\c_1\\c_2\\c_3\\c_4\end{bmatrix}=
\begin{bmatrix}0\\1\\0\\1\\2\end{bmatrix}
\]

\begin{enumerate}
\def\labelenumi{(\alph{enumi})}
\setcounter{enumi}{2}
\tightlist
\item
  code
\end{enumerate}

    \begin{tcolorbox}[breakable, size=fbox, boxrule=1pt, pad at break*=1mm,colback=cellbackground, colframe=cellborder]
\prompt{In}{incolor}{206}{\boxspacing}
\begin{Verbatim}[commandchars=\\\{\}]
\PY{k}{local}\PY{+w}{ }\PY{n}{x}\PY{p}{,}\PY{+w}{ }\PY{n}{M}\PY{p}{,}\PY{+w}{ }\PY{n}{y}
\PY{k}{let}\PY{+w}{ }\PY{n}{x}\PY{p}{,}\PY{+w}{ }\PY{n}{M}\PY{p}{,}\PY{+w}{ }\PY{n}{y}
\PY{+w}{    }\PY{n}{x}\PY{+w}{ }\PY{o}{=}\PY{+w}{ }\PY{p}{[}\PY{o}{\PYZhy{}}\PY{l+m+mi}{1}\PY{p}{,}\PY{+w}{ }\PY{l+m+mi}{0}\PY{p}{,}\PY{+w}{ }\PY{l+m+mi}{2}\PY{p}{,}\PY{+w}{ }\PY{l+m+mi}{3}\PY{p}{,}\PY{+w}{ }\PY{l+m+mi}{4}\PY{p}{]}
\PY{+w}{    }\PY{n}{M}\PY{+w}{ }\PY{o}{=}\PY{+w}{ }\PY{p}{[}\PY{+w}{ }\PY{n}{x}\PY{p}{[}\PY{n}{i}\PY{p}{]}\PY{+w}{ }\PY{o}{\PYZca{}}\PY{+w}{ }\PY{n}{j}\PY{+w}{ }\PY{k}{for}\PY{+w}{ }\PY{n}{i}\PY{+w}{ }\PY{k}{in}\PY{+w}{ }\PY{l+m+mi}{1}\PY{+w}{ }\PY{o}{:}\PY{+w}{ }\PY{n}{length}\PY{p}{(}\PY{n}{x}\PY{p}{)}\PY{p}{,}\PY{+w}{ }\PY{n}{j}\PY{+w}{ }\PY{k}{in}\PY{+w}{ }\PY{l+m+mi}{0}\PY{o}{:}\PY{n}{length}\PY{p}{(}\PY{n}{x}\PY{p}{)}\PY{+w}{ }\PY{o}{\PYZhy{}}\PY{+w}{ }\PY{l+m+mi}{1}\PY{+w}{ }\PY{p}{]}
\PY{+w}{    }\PY{n}{y}\PY{+w}{ }\PY{o}{=}\PY{+w}{ }\PY{p}{[}\PY{l+m+mi}{0}\PY{p}{;}\PY{+w}{ }\PY{l+m+mi}{1}\PY{p}{;}\PY{+w}{ }\PY{l+m+mi}{0}\PY{p}{;}\PY{+w}{ }\PY{l+m+mi}{1}\PY{p}{;}\PY{+w}{ }\PY{l+m+mi}{2}\PY{p}{]}

\PY{+w}{    }\PY{n}{println}\PY{p}{(}\PY{n}{M}\PY{p}{)}

\PY{+w}{    }\PY{n}{println}\PY{p}{(}\PY{l+s}{\PYZdq{}}\PY{l+s}{c: }\PY{l+s}{\PYZdq{}}\PY{p}{,}\PY{+w}{ }\PY{n}{M}\PY{+w}{ }\PY{o}{\PYZbs{}}\PY{+w}{ }\PY{n}{y}\PY{p}{)}
\PY{k}{end}
\end{Verbatim}
\end{tcolorbox}

    \begin{Verbatim}[commandchars=\\\{\}]
[1 -1 1 -1 1; 1 0 0 0 0; 1 2 4 8 16; 1 3 9 27 81; 1 4 16 64 256]
c: [1.0000000000000004, -0.44999999999999896, -0.8249999999999996,
0.5499999999999998, -0.07499999999999998]
    \end{Verbatim}

    Hence the polynomial is \(p(t)=1-0.45t-0.825t^2+0.55t^3-0.075t^4\)

    \hypertarget{section-2.1-3}{%
\subsection{Section 2.1 \#3}\label{section-2.1-3}}

Hint: use the horner function from the textbook (Function 1.3.2) or
Polynomial from the package Polynomials to create and evaluate a
polynomial from its coefficients.

    \begin{tcolorbox}[breakable, size=fbox, boxrule=1pt, pad at break*=1mm,colback=cellbackground, colframe=cellborder]
\prompt{In}{incolor}{207}{\boxspacing}
\begin{Verbatim}[commandchars=\\\{\}]
\PY{k}{local}\PY{+w}{ }\PY{n}{x0}\PY{p}{,}\PY{+w}{ }\PY{n}{x}\PY{p}{,}\PY{+w}{ }\PY{n}{p\PYZus{}USA}\PY{p}{,}\PY{+w}{ }\PY{n}{p\PYZus{}India}\PY{p}{,}\PY{+w}{ }\PY{n}{p\PYZus{}Poland}

\PY{n}{x0}\PY{+w}{ }\PY{o}{=}\PY{+w}{ }\PY{p}{[}\PY{l+m+mi}{1990}\PY{p}{,}\PY{+w}{ }\PY{l+m+mi}{2000}\PY{p}{,}\PY{+w}{ }\PY{l+m+mi}{2010}\PY{p}{,}\PY{+w}{ }\PY{l+m+mi}{2020}\PY{p}{]}\PY{p}{;}
\PY{n}{x}\PY{+w}{ }\PY{o}{=}\PY{+w}{ }\PY{n}{x0}\PY{+w}{ }\PY{o}{.\PYZhy{}}\PY{+w}{ }\PY{l+m+mi}{1990}\PY{p}{;}
\PY{n}{p\PYZus{}USA}\PY{+w}{ }\PY{o}{=}\PY{+w}{ }\PY{p}{[}\PY{l+m+mf}{252.120}\PY{p}{,}\PY{+w}{ }\PY{l+m+mf}{281.711}\PY{p}{,}\PY{+w}{ }\PY{l+m+mf}{309.011}\PY{p}{,}\PY{+w}{ }\PY{l+m+mf}{331.003}\PY{p}{]}\PY{p}{;}
\PY{n}{p\PYZus{}India}\PY{+w}{ }\PY{o}{=}\PY{+w}{ }\PY{p}{[}\PY{l+m+mf}{873.278}\PY{p}{,}\PY{+w}{ }\PY{l+m+mf}{1056.576}\PY{p}{,}\PY{+w}{ }\PY{l+m+mf}{1234.281}\PY{p}{,}\PY{+w}{ }\PY{l+m+mf}{1\PYZus{}380.004}\PY{p}{]}\PY{p}{;}
\PY{n}{p\PYZus{}Poland}\PY{+w}{ }\PY{o}{=}\PY{+w}{ }\PY{p}{[}\PY{l+m+mf}{37.960}\PY{p}{,}\PY{+w}{ }\PY{l+m+mf}{38.557}\PY{p}{,}\PY{+w}{ }\PY{l+m+mf}{38.330}\PY{p}{,}\PY{+w}{ }\PY{l+m+mf}{37.847}\PY{p}{]}\PY{p}{;}
\end{Verbatim}
\end{tcolorbox}

    \begin{enumerate}
\def\labelenumi{(\alph{enumi})}
\tightlist
\item
  The population of the USA in 2005 is estimated to be 295.84 millions.
\end{enumerate}

    \begin{tcolorbox}[breakable, size=fbox, boxrule=1pt, pad at break*=1mm,colback=cellbackground, colframe=cellborder]
\prompt{In}{incolor}{208}{\boxspacing}
\begin{Verbatim}[commandchars=\\\{\}]
\PY{n}{M}\PY{+w}{ }\PY{o}{=}\PY{+w}{ }\PY{p}{[}\PY{+w}{ }\PY{n}{x}\PY{p}{[}\PY{n}{i}\PY{p}{]}\PY{+w}{ }\PY{o}{\PYZca{}}\PY{+w}{ }\PY{n}{j}\PY{+w}{ }\PY{k}{for}\PY{+w}{ }\PY{n}{i}\PY{+w}{ }\PY{k}{in}\PY{+w}{ }\PY{l+m+mi}{1}\PY{+w}{ }\PY{o}{:}\PY{+w}{ }\PY{n}{length}\PY{p}{(}\PY{n}{x}\PY{p}{)}\PY{p}{,}\PY{+w}{ }\PY{n}{j}\PY{+w}{ }\PY{k}{in}\PY{+w}{ }\PY{l+m+mi}{0}\PY{o}{:}\PY{n}{length}\PY{p}{(}\PY{n}{x}\PY{p}{)}\PY{+w}{ }\PY{o}{\PYZhy{}}\PY{+w}{ }\PY{l+m+mi}{1}\PY{+w}{ }\PY{p}{]}
\PY{n}{c\PYZus{}USA}\PY{+w}{ }\PY{o}{=}\PY{+w}{ }\PY{n}{M}\PY{+w}{ }\PY{o}{\PYZbs{}}\PY{+w}{ }\PY{n}{p\PYZus{}USA}

\PY{n}{f\PYZus{}USA}\PY{p}{(}\PY{n}{x}\PY{p}{)}\PY{+w}{ }\PY{o}{=}\PY{+w}{ }\PY{n}{sum}\PY{p}{(}\PY{n}{c\PYZus{}USA}\PY{+w}{ }\PY{o}{.*}\PY{+w}{ }\PY{p}{(}\PY{n}{x}\PY{+w}{ }\PY{o}{.\PYZca{}}\PY{+w}{ }\PY{p}{[}\PY{l+m+mi}{0}\PY{p}{,}\PY{+w}{ }\PY{l+m+mi}{1}\PY{p}{,}\PY{+w}{ }\PY{l+m+mi}{2}\PY{p}{,}\PY{+w}{ }\PY{l+m+mi}{3}\PY{p}{]}\PY{p}{)}\PY{p}{)}

\PY{n}{f\PYZus{}USA}\PY{p}{(}\PY{l+m+mi}{2005}\PY{+w}{ }\PY{o}{\PYZhy{}}\PY{+w}{ }\PY{l+m+mi}{1990}\PY{p}{)}
\end{Verbatim}
\end{tcolorbox}

    
    \begin{Verbatim}[commandchars=\\\{\}]
295.83593750000006
    \end{Verbatim}

    
    \begin{enumerate}
\def\labelenumi{(\alph{enumi})}
\setcounter{enumi}{1}
\tightlist
\item
\end{enumerate}

    \begin{tcolorbox}[breakable, size=fbox, boxrule=1pt, pad at break*=1mm,colback=cellbackground, colframe=cellborder]
\prompt{In}{incolor}{209}{\boxspacing}
\begin{Verbatim}[commandchars=\\\{\}]
\PY{k}{using}\PY{+w}{ }\PY{n}{Plots}
\PY{k}{using}\PY{+w}{ }\PY{n}{Roots}

\PY{n}{c\PYZus{}Poland}\PY{o}{=}\PY{+w}{ }\PY{n}{M}\PY{+w}{ }\PY{o}{\PYZbs{}}\PY{+w}{ }\PY{n}{p\PYZus{}Poland}

\PY{n}{f\PYZus{}Poland}\PY{p}{(}\PY{n}{x}\PY{p}{)}\PY{+w}{ }\PY{o}{=}\PY{+w}{ }\PY{n}{sum}\PY{p}{(}\PY{n}{c\PYZus{}Poland}\PY{+w}{ }\PY{o}{.*}\PY{+w}{ }\PY{p}{(}\PY{p}{(}\PY{n}{x}\PY{+w}{ }\PY{o}{\PYZhy{}}\PY{+w}{ }\PY{l+m+mi}{1990}\PY{p}{)}\PY{+w}{ }\PY{o}{.\PYZca{}}\PY{+w}{ }\PY{p}{[}\PY{l+m+mi}{0}\PY{p}{,}\PY{+w}{ }\PY{l+m+mi}{1}\PY{p}{,}\PY{+w}{ }\PY{l+m+mi}{2}\PY{p}{,}\PY{+w}{ }\PY{l+m+mi}{3}\PY{p}{]}\PY{p}{)}\PY{p}{)}

\PY{n}{plot}\PY{p}{(}\PY{n}{f\PYZus{}Poland}\PY{p}{,}\PY{+w}{ }\PY{l+m+mi}{1990}\PY{p}{,}\PY{+w}{ }\PY{l+m+mi}{2020}\PY{p}{,}\PY{+w}{ }
\PY{+w}{    }\PY{n}{title}\PY{o}{=}\PY{l+s}{\PYZdq{}}\PY{l+s}{Graph of the population of Poland}\PY{l+s}{\PYZdq{}}\PY{p}{,}\PY{+w}{ }
\PY{+w}{    }\PY{n}{label}\PY{o}{=}\PY{l+s}{\PYZdq{}}\PY{l+s}{Population}\PY{l+s}{\PYZdq{}}\PY{p}{,}
\PY{+w}{    }\PY{n}{linewidth}\PY{o}{=}\PY{l+m+mi}{3}\PY{p}{,}
\PY{+w}{    }\PY{n}{xlabel}\PY{o}{=}\PY{l+s}{\PYZdq{}}\PY{l+s}{Year}\PY{l+s}{\PYZdq{}}\PY{p}{,}
\PY{+w}{    }\PY{n}{ylabel}\PY{o}{=}\PY{l+s}{\PYZdq{}}\PY{l+s}{Population (million)}\PY{l+s}{\PYZdq{}}\PY{p}{,}
\PY{+w}{    }\PY{c}{\PYZsh{} labelfont=font(14, \PYZdq{}Fira Code\PYZdq{})}
\PY{+w}{    }\PY{p}{)}
\end{Verbatim}
\end{tcolorbox}

    \begin{center}
    \adjustimage{max size={0.9\linewidth}{0.9\paperheight}}{hw-02_files/hw-02_9_0.pdf}
    \end{center}
    { \hspace*{\fill} \\}
    
    From the graph we can see that there is only one peak within
\([1990, 2020]\), likely to be 2000, hence we only need to find the zero
point of the first order derivative of the interpolating function
\texttt{fp\_Poland}.

It's 1st order derivative is \(p'(t)=c_1+2c_2t+3c_3t^2\).

    \begin{tcolorbox}[breakable, size=fbox, boxrule=1pt, pad at break*=1mm,colback=cellbackground, colframe=cellborder]
\prompt{In}{incolor}{210}{\boxspacing}
\begin{Verbatim}[commandchars=\\\{\}]
\PY{c}{\PYZsh{} derivative of the interpolating function}
\PY{n}{fp\PYZus{}Poland}\PY{p}{(}\PY{n}{x}\PY{p}{)}\PY{+w}{ }\PY{o}{=}\PY{+w}{ }\PY{n}{sum}\PY{p}{(}\PY{n}{c\PYZus{}Poland}\PY{+w}{ }\PY{o}{.*}\PY{+w}{ }\PY{p}{[}\PY{l+m+mi}{0}\PY{p}{,}\PY{+w}{ }\PY{l+m+mi}{1}\PY{p}{,}\PY{+w}{ }\PY{l+m+mi}{2}\PY{p}{,}\PY{+w}{ }\PY{l+m+mi}{3}\PY{p}{]}\PY{+w}{ }\PY{o}{.*}\PY{+w}{ }\PY{p}{(}\PY{p}{(}\PY{n}{x}\PY{+w}{ }\PY{o}{\PYZhy{}}\PY{+w}{ }\PY{l+m+mi}{1990}\PY{p}{)}\PY{+w}{ }\PY{o}{.\PYZca{}}\PY{+w}{ }\PY{p}{[}\PY{l+m+mi}{0}\PY{p}{,}\PY{+w}{ }\PY{l+m+mi}{0}\PY{p}{,}\PY{+w}{ }\PY{l+m+mi}{1}\PY{p}{,}\PY{+w}{ }\PY{l+m+mi}{2}\PY{p}{]}\PY{p}{)}\PY{p}{)}
\PY{n}{zero}\PY{+w}{ }\PY{o}{=}\PY{+w}{ }\PY{n}{find\PYZus{}zero}\PY{p}{(}\PY{n}{fp\PYZus{}Poland}\PY{p}{,}\PY{+w}{ }\PY{l+m+mi}{2000}\PY{p}{)}

\PY{n}{println}\PY{p}{(}\PY{l+s}{\PYZdq{}}\PY{l+s}{Peak year: }\PY{l+s}{\PYZdq{}}\PY{p}{,}\PY{+w}{ }\PY{n}{zero}\PY{p}{)}
\PY{n}{println}\PY{p}{(}\PY{l+s}{\PYZdq{}}\PY{l+s}{Peak population (million): }\PY{l+s}{\PYZdq{}}\PY{p}{,}\PY{+w}{ }\PY{n}{f\PYZus{}Poland}\PY{p}{(}\PY{n}{round}\PY{p}{(}\PY{n}{zero}\PY{p}{)}\PY{p}{)}\PY{p}{)}
\end{Verbatim}
\end{tcolorbox}

    \begin{Verbatim}[commandchars=\\\{\}]
Peak year: 2001.1411616909593
Peak population (million): 38.562008000000006
    \end{Verbatim}

    Therefore the peak locates at year 2001, with maximum population about
38.56 million.

    \begin{enumerate}
\def\labelenumi{(\alph{enumi})}
\setcounter{enumi}{2}
\tightlist
\item
  plots:
\end{enumerate}

    \begin{tcolorbox}[breakable, size=fbox, boxrule=1pt, pad at break*=1mm,colback=cellbackground, colframe=cellborder]
\prompt{In}{incolor}{211}{\boxspacing}
\begin{Verbatim}[commandchars=\\\{\}]
\PY{n}{c\PYZus{}India}\PY{o}{=}\PY{+w}{ }\PY{n}{M}\PY{+w}{ }\PY{o}{\PYZbs{}}\PY{+w}{ }\PY{n}{p\PYZus{}India}

\PY{n}{f\PYZus{}India}\PY{p}{(}\PY{n}{x}\PY{p}{)}\PY{+w}{ }\PY{o}{=}\PY{+w}{ }\PY{n}{sum}\PY{p}{(}\PY{n}{c\PYZus{}India}\PY{+w}{ }\PY{o}{.*}\PY{+w}{ }\PY{p}{(}\PY{p}{(}\PY{n}{x}\PY{+w}{ }\PY{o}{\PYZhy{}}\PY{+w}{ }\PY{l+m+mi}{1990}\PY{p}{)}\PY{+w}{ }\PY{o}{.\PYZca{}}\PY{+w}{ }\PY{p}{[}\PY{l+m+mi}{0}\PY{p}{,}\PY{+w}{ }\PY{l+m+mi}{1}\PY{p}{,}\PY{+w}{ }\PY{l+m+mi}{2}\PY{p}{,}\PY{+w}{ }\PY{l+m+mi}{3}\PY{p}{]}\PY{p}{)}\PY{p}{)}

\PY{n}{plot}\PY{p}{(}\PY{n}{f\PYZus{}India}\PY{p}{,}\PY{+w}{ }\PY{l+m+mi}{1985}\PY{p}{,}\PY{+w}{ }\PY{l+m+mi}{2025}\PY{p}{,}\PY{+w}{ }
\PY{+w}{    }\PY{n}{title}\PY{o}{=}\PY{l+s}{\PYZdq{}}\PY{l+s}{Graph of the population of India}\PY{l+s}{\PYZdq{}}\PY{p}{,}\PY{+w}{ }
\PY{+w}{    }\PY{n}{label}\PY{o}{=}\PY{l+s}{\PYZdq{}}\PY{l+s}{Population}\PY{l+s}{\PYZdq{}}\PY{p}{,}
\PY{+w}{    }\PY{n}{linewidth}\PY{o}{=}\PY{l+m+mi}{3}\PY{p}{,}
\PY{+w}{    }\PY{n}{xlabel}\PY{o}{=}\PY{l+s}{\PYZdq{}}\PY{l+s}{Year}\PY{l+s}{\PYZdq{}}\PY{p}{,}
\PY{+w}{    }\PY{n}{ylabel}\PY{o}{=}\PY{l+s}{\PYZdq{}}\PY{l+s}{Population (million)}\PY{l+s}{\PYZdq{}}
\PY{+w}{    }\PY{c}{\PYZsh{} labelfont=font(14, \PYZdq{}Fira Code\PYZdq{})}
\PY{+w}{    }\PY{p}{)}
\PY{n}{scatter!}\PY{p}{(}\PY{n}{x0}\PY{p}{,}\PY{+w}{ }\PY{n}{p\PYZus{}India}\PY{p}{,}\PY{+w}{ }\PY{n}{label}\PY{o}{=}\PY{l+s}{\PYZdq{}}\PY{l+s}{actual}\PY{l+s}{\PYZdq{}}\PY{p}{,}\PY{+w}{ }\PY{n}{legend}\PY{o}{=}\PY{l+s+ss}{:topleft}\PY{p}{)}
\end{Verbatim}
\end{tcolorbox}

    \begin{center}
    \adjustimage{max size={0.9\linewidth}{0.9\paperheight}}{hw-02_files/hw-02_14_0.pdf}
    \end{center}
    { \hspace*{\fill} \\}
    
    \hypertarget{section-2.3-7}{%
\subsection{Section 2.3 \#7}\label{section-2.3-7}}

Demo 2.3.6 showed solutions of \(Ax=b\), where \(A\) and \(b\) are
computed below as a function of \(\alpha\) and \(\beta\). Use Function
2.3.5 to solve with \(\alpha = 0.1\) and \(\beta = 10^i\),
\(i = 1, \dots, 12\), tabulating the values of \(\beta\) and
\(|x_1-1|\). (This kind of behavior is explained in Section 2.8.)

The code below will help you create the table and display it nicely.

    \begin{tcolorbox}[breakable, size=fbox, boxrule=1pt, pad at break*=1mm,colback=cellbackground, colframe=cellborder]
\prompt{In}{incolor}{212}{\boxspacing}
\begin{Verbatim}[commandchars=\\\{\}]
\PY{c}{\PYZsh{} you might want the textbook\PYZsq{}s backsub function too (Function 2.3.5)}
\PY{k}{function}\PY{+w}{ }\PY{n}{backsub}\PY{p}{(}\PY{n}{U}\PY{p}{,}\PY{n}{b}\PY{p}{)}
\PY{+w}{    }\PY{n}{n}\PY{+w}{ }\PY{o}{=}\PY{+w}{ }\PY{n}{size}\PY{p}{(}\PY{n}{U}\PY{p}{,}\PY{l+m+mi}{1}\PY{p}{)}
\PY{+w}{    }\PY{n}{x}\PY{+w}{ }\PY{o}{=}\PY{+w}{ }\PY{n}{zeros}\PY{p}{(}\PY{n}{n}\PY{p}{)}
\PY{+w}{    }\PY{n}{x}\PY{p}{[}\PY{n}{n}\PY{p}{]}\PY{+w}{ }\PY{o}{=}\PY{+w}{ }\PY{n}{b}\PY{p}{[}\PY{n}{n}\PY{p}{]}\PY{o}{/}\PY{n}{U}\PY{p}{[}\PY{n}{n}\PY{p}{,}\PY{n}{n}\PY{p}{]}
\PY{+w}{    }\PY{k}{for}\PY{+w}{ }\PY{n}{i}\PY{+w}{ }\PY{k}{in}\PY{+w}{ }\PY{n}{n}\PY{o}{\PYZhy{}}\PY{l+m+mi}{1}\PY{o}{:}\PY{o}{\PYZhy{}}\PY{l+m+mi}{1}\PY{o}{:}\PY{l+m+mi}{1}
\PY{+w}{        }\PY{n}{s}\PY{+w}{ }\PY{o}{=}\PY{+w}{ }\PY{n}{sum}\PY{p}{(}\PY{+w}{ }\PY{n}{U}\PY{p}{[}\PY{n}{i}\PY{p}{,}\PY{n}{j}\PY{p}{]}\PY{o}{*}\PY{n}{x}\PY{p}{[}\PY{n}{j}\PY{p}{]}\PY{+w}{ }\PY{k}{for}\PY{+w}{ }\PY{n}{j}\PY{+w}{ }\PY{k}{in}\PY{+w}{ }\PY{n}{i}\PY{o}{+}\PY{l+m+mi}{1}\PY{o}{:}\PY{n}{n}\PY{+w}{ }\PY{p}{)}
\PY{+w}{        }\PY{n}{x}\PY{p}{[}\PY{n}{i}\PY{p}{]}\PY{+w}{ }\PY{o}{=}\PY{+w}{ }\PY{p}{(}\PY{+w}{ }\PY{n}{b}\PY{p}{[}\PY{n}{i}\PY{p}{]}\PY{+w}{ }\PY{o}{\PYZhy{}}\PY{+w}{ }\PY{n}{s}\PY{+w}{ }\PY{p}{)}\PY{+w}{ }\PY{o}{/}\PY{+w}{ }\PY{n}{U}\PY{p}{[}\PY{n}{i}\PY{p}{,}\PY{n}{i}\PY{p}{]}
\PY{+w}{    }\PY{k}{end}
\PY{+w}{    }\PY{k}{return}\PY{+w}{ }\PY{n}{x}
\PY{k}{end}
\end{Verbatim}
\end{tcolorbox}

    
    \begin{Verbatim}[commandchars=\\\{\}]
backsub (generic function with 1 method)
    \end{Verbatim}

    
    \begin{tcolorbox}[breakable, size=fbox, boxrule=1pt, pad at break*=1mm,colback=cellbackground, colframe=cellborder]
\prompt{In}{incolor}{213}{\boxspacing}
\begin{Verbatim}[commandchars=\\\{\}]
\PY{k}{using}\PY{+w}{ }\PY{n}{PrettyTables}
\PY{k}{using}\PY{+w}{ }\PY{n}{LinearAlgebra}

\PY{k}{local}\PY{+w}{ }\PY{n}{a}\PY{p}{,}\PY{+w}{ }\PY{n}{bs}\PY{p}{,}\PY{+w}{ }\PY{n}{errs}
\PY{k}{let}\PY{+w}{ }\PY{n}{a}\PY{p}{,}\PY{+w}{ }\PY{n}{bs}\PY{p}{,}\PY{+w}{ }\PY{n}{errs}
\PY{+w}{    }\PY{n}{a}\PY{+w}{ }\PY{o}{=}\PY{+w}{ }\PY{l+m+mf}{0.1}
\PY{+w}{    }\PY{n}{bs}\PY{+w}{ }\PY{o}{=}\PY{+w}{ }\PY{l+m+mi}{10}\PY{+w}{ }\PY{o}{.\PYZca{}}\PY{+w}{ }\PY{p}{(}\PY{l+m+mi}{1}\PY{+w}{ }\PY{o}{:}\PY{+w}{ }\PY{l+m+mi}{12}\PY{p}{)}

\PY{+w}{    }\PY{n}{fA}\PY{p}{(}\PY{n}{a}\PY{p}{,}\PY{+w}{ }\PY{n}{b}\PY{p}{)}\PY{+w}{ }\PY{o}{=}\PY{+w}{ }\PY{n}{diagm}\PY{p}{(}\PY{l+m+mi}{0}\PY{+w}{ }\PY{o}{=\PYZgt{}}\PY{+w}{ }\PY{n}{ones}\PY{p}{(}\PY{l+m+mi}{5}\PY{p}{)}\PY{p}{,}\PY{+w}{ }\PY{l+m+mi}{1}\PY{+w}{ }\PY{o}{=\PYZgt{}}\PY{+w}{ }\PY{p}{[}\PY{o}{\PYZhy{}}\PY{l+m+mi}{1}\PY{p}{,}\PY{o}{\PYZhy{}}\PY{l+m+mi}{1}\PY{p}{,}\PY{o}{\PYZhy{}}\PY{l+m+mi}{1}\PY{p}{,}\PY{o}{\PYZhy{}}\PY{l+m+mi}{1}\PY{p}{]}\PY{p}{,}\PY{+w}{ }\PY{l+m+mi}{3}\PY{+w}{ }\PY{o}{=\PYZgt{}}\PY{+w}{ }\PY{p}{[}\PY{n}{a}\PY{+w}{ }\PY{o}{\PYZhy{}}\PY{+w}{ }\PY{n}{b}\PY{p}{,}\PY{+w}{ }\PY{l+m+mi}{0}\PY{p}{]}\PY{p}{,}\PY{+w}{ }\PY{l+m+mi}{4}\PY{+w}{ }\PY{o}{=\PYZgt{}}\PY{+w}{ }\PY{p}{[}\PY{n}{b}\PY{p}{]}\PY{p}{)}
\PY{+w}{    }\PY{n}{fb}\PY{p}{(}\PY{n}{a}\PY{p}{)}\PY{+w}{ }\PY{o}{=}\PY{+w}{ }\PY{p}{[}\PY{n}{a}\PY{p}{;}\PY{+w}{ }\PY{l+m+mi}{0}\PY{p}{;}\PY{+w}{ }\PY{l+m+mi}{0}\PY{p}{;}\PY{+w}{ }\PY{l+m+mi}{0}\PY{p}{;}\PY{+w}{ }\PY{l+m+mi}{1}\PY{p}{]}

\PY{+w}{    }\PY{n}{backsub}\PY{p}{(}\PY{n}{fA}\PY{p}{(}\PY{n}{a}\PY{p}{,}\PY{+w}{ }\PY{n}{bs}\PY{p}{[}\PY{l+m+mi}{10}\PY{p}{]}\PY{p}{)}\PY{p}{,}\PY{+w}{ }\PY{n}{fb}\PY{p}{(}\PY{n}{a}\PY{p}{)}\PY{p}{)}

\PY{+w}{    }\PY{n}{errs}\PY{+w}{ }\PY{o}{=}\PY{+w}{ }\PY{p}{[}\PY{n}{abs}\PY{p}{(}\PY{n}{backsub}\PY{p}{(}\PY{n}{fA}\PY{p}{(}\PY{n}{a}\PY{p}{,}\PY{+w}{ }\PY{n}{b}\PY{p}{)}\PY{p}{,}\PY{+w}{ }\PY{n}{fb}\PY{p}{(}\PY{n}{a}\PY{p}{)}\PY{p}{)}\PY{p}{[}\PY{l+m+mi}{1}\PY{p}{]}\PY{+w}{ }\PY{o}{\PYZhy{}}\PY{+w}{ }\PY{l+m+mi}{1}\PY{p}{)}\PY{+w}{ }\PY{k}{for}\PY{+w}{ }\PY{n}{b}\PY{+w}{ }\PY{k}{in}\PY{+w}{ }\PY{n}{bs}\PY{p}{]}

\PY{+w}{    }\PY{n}{pretty\PYZus{}table}\PY{p}{(}\PY{p}{[}\PY{n}{bs}\PY{+w}{ }\PY{n}{errs}\PY{p}{]}\PY{p}{,}\PY{+w}{ }\PY{n}{crop}\PY{+w}{ }\PY{o}{=}\PY{+w}{ }\PY{l+s+ss}{:none}\PY{p}{,}\PY{+w}{ }\PY{n}{header}\PY{+w}{ }\PY{o}{=}\PY{+w}{ }\PY{p}{[}\PY{l+s}{\PYZdq{}}\PY{l+s}{β}\PY{l+s}{\PYZdq{}}\PY{p}{,}\PY{+w}{ }\PY{l+s}{\PYZdq{}}\PY{l+s}{|x₁\PYZhy{}1|}\PY{l+s}{\PYZdq{}}\PY{p}{]}\PY{p}{)}
\PY{k}{end}
\end{Verbatim}
\end{tcolorbox}

    \begin{Verbatim}[commandchars=\\\{\}]
┌──────────┬─────────────┐
│\textbf{        β }│\textbf{      |x₁-1| }│
├──────────┼─────────────┤
│     10.0 │ 4.44089e-16 │
│    100.0 │ 5.77316e-15 │
│   1000.0 │ 2.27596e-14 │
│  10000.0 │  3.6382e-13 │
│ 100000.0 │ 5.82079e-12 │
│    1.0e6 │ 2.32832e-11 │
│    1.0e7 │ 3.72529e-10 │
│    1.0e8 │  5.96046e-9 │
│    1.0e9 │  2.38419e-8 │
│   1.0e10 │   3.8147e-7 │
│   1.0e11 │  6.10352e-6 │
│   1.0e12 │  2.44141e-5 │
└──────────┴─────────────┘
    \end{Verbatim}

    As \(\alpha\) increases exponentially, the error of the solution \(x\)
increases.

    \hypertarget{section-2.4-6}{%
\subsection{Section 2.4 \#6}\label{section-2.4-6}}

The \texttt{det} function is in LinearAlgebra. You may want Function
2.4.6 from the book to help too.

The function \texttt{diag} can be used to extract the diagonal elements
of a matrix into a vector. The function \texttt{prod} is like
\texttt{sum}, i.e., instead of adding the entries of a vector, it
computes the products of the entries.

    \textbf{Answer:}

\begin{enumerate}
\def\labelenumi{(\alph{enumi})}
\item
  If \(A=LU\), by the definition and theorems of determinant,
  \[|L|=\prod_{i=1}^n L_{ii}=1,\quad |U|=\prod_{i=1}^n U_{ii}, \quad |A|=|LU|=|L||U|=\prod_{i=1}^n U_{ii}\]
\item
\end{enumerate}

    \begin{tcolorbox}[breakable, size=fbox, boxrule=1pt, pad at break*=1mm,colback=cellbackground, colframe=cellborder]
\prompt{In}{incolor}{214}{\boxspacing}
\begin{Verbatim}[commandchars=\\\{\}]
\PY{k}{function}\PY{+w}{ }\PY{n}{lufact}\PY{p}{(}\PY{n}{A}\PY{p}{)}
\PY{+w}{    }\PY{n}{n}\PY{+w}{ }\PY{o}{=}\PY{+w}{ }\PY{n}{size}\PY{p}{(}\PY{n}{A}\PY{p}{,}\PY{l+m+mi}{1}\PY{p}{)}
\PY{+w}{    }\PY{n}{L}\PY{+w}{ }\PY{o}{=}\PY{+w}{ }\PY{n}{diagm}\PY{p}{(}\PY{n}{ones}\PY{p}{(}\PY{n}{n}\PY{p}{)}\PY{p}{)}\PY{+w}{   }\PY{c}{\PYZsh{} ones on main diagonal, zeros elsewhere}
\PY{+w}{    }\PY{n}{U}\PY{+w}{ }\PY{o}{=}\PY{+w}{ }\PY{n}{zeros}\PY{p}{(}\PY{n}{n}\PY{p}{,}\PY{n}{n}\PY{p}{)}
\PY{+w}{    }\PY{n}{Aₖ}\PY{+w}{ }\PY{o}{=}\PY{+w}{ }\PY{n}{float}\PY{p}{(}\PY{n}{copy}\PY{p}{(}\PY{n}{A}\PY{p}{)}\PY{p}{)}

\PY{+w}{    }\PY{c}{\PYZsh{} Reduction by outer products}
\PY{+w}{    }\PY{k}{for}\PY{+w}{ }\PY{n}{k}\PY{+w}{ }\PY{k}{in}\PY{+w}{ }\PY{l+m+mi}{1}\PY{o}{:}\PY{n}{n}\PY{o}{\PYZhy{}}\PY{l+m+mi}{1}
\PY{+w}{        }\PY{n}{U}\PY{p}{[}\PY{n}{k}\PY{p}{,}\PY{o}{:}\PY{p}{]}\PY{+w}{ }\PY{o}{=}\PY{+w}{ }\PY{n}{Aₖ}\PY{p}{[}\PY{n}{k}\PY{p}{,}\PY{o}{:}\PY{p}{]}
\PY{+w}{        }\PY{n}{L}\PY{p}{[}\PY{o}{:}\PY{p}{,}\PY{n}{k}\PY{p}{]}\PY{+w}{ }\PY{o}{=}\PY{+w}{ }\PY{n}{Aₖ}\PY{p}{[}\PY{o}{:}\PY{p}{,}\PY{n}{k}\PY{p}{]}\PY{o}{/}\PY{n}{U}\PY{p}{[}\PY{n}{k}\PY{p}{,}\PY{n}{k}\PY{p}{]}
\PY{+w}{        }\PY{n}{Aₖ}\PY{+w}{ }\PY{o}{\PYZhy{}=}\PY{+w}{ }\PY{n}{L}\PY{p}{[}\PY{o}{:}\PY{p}{,}\PY{n}{k}\PY{p}{]}\PY{o}{*}\PY{n}{U}\PY{p}{[}\PY{n}{k}\PY{p}{,}\PY{o}{:}\PY{p}{]}\PY{o}{\PYZsq{}}
\PY{+w}{    }\PY{k}{end}
\PY{+w}{    }\PY{n}{U}\PY{p}{[}\PY{n}{n}\PY{p}{,}\PY{n}{n}\PY{p}{]}\PY{+w}{ }\PY{o}{=}\PY{+w}{ }\PY{n}{Aₖ}\PY{p}{[}\PY{n}{n}\PY{p}{,}\PY{n}{n}\PY{p}{]}
\PY{+w}{    }\PY{k}{return}\PY{+w}{ }\PY{n}{LowerTriangular}\PY{p}{(}\PY{n}{L}\PY{p}{)}\PY{p}{,}\PY{n}{UpperTriangular}\PY{p}{(}\PY{n}{U}\PY{p}{)}
\PY{k}{end}
\end{Verbatim}
\end{tcolorbox}

    
    \begin{Verbatim}[commandchars=\\\{\}]
lufact (generic function with 1 method)
    \end{Verbatim}

    
    \begin{tcolorbox}[breakable, size=fbox, boxrule=1pt, pad at break*=1mm,colback=cellbackground, colframe=cellborder]
\prompt{In}{incolor}{215}{\boxspacing}
\begin{Verbatim}[commandchars=\\\{\}]
\PY{k}{function}\PY{+w}{ }\PY{n}{determinant}\PY{p}{(}\PY{n}{A}\PY{p}{)}
\PY{+w}{    }\PY{n}{\PYZus{}}\PY{p}{,}\PY{+w}{ }\PY{n}{U}\PY{+w}{ }\PY{o}{=}\PY{+w}{ }\PY{n}{lufact}\PY{p}{(}\PY{n}{A}\PY{p}{)}
\PY{+w}{    }\PY{k}{return}\PY{+w}{ }\PY{n}{prod}\PY{p}{(}\PY{n}{diag}\PY{p}{(}\PY{n}{U}\PY{p}{)}\PY{p}{)}
\PY{k}{end}
\end{Verbatim}
\end{tcolorbox}

    
    \begin{Verbatim}[commandchars=\\\{\}]
determinant (generic function with 1 method)
    \end{Verbatim}

    
    \begin{tcolorbox}[breakable, size=fbox, boxrule=1pt, pad at break*=1mm,colback=cellbackground, colframe=cellborder]
\prompt{In}{incolor}{216}{\boxspacing}
\begin{Verbatim}[commandchars=\\\{\}]
\PY{k}{local}\PY{+w}{ }\PY{n}{A1}\PY{p}{,}\PY{+w}{ }\PY{n}{A2}
\PY{k}{let}\PY{+w}{ }\PY{n}{A1}\PY{p}{,}\PY{+w}{ }\PY{n}{A2}

\PY{+w}{   }\PY{n}{A1}\PY{+w}{ }\PY{o}{=}\PY{+w}{ }\PY{p}{[}\PY{+w}{              }
\PY{+w}{         }\PY{l+m+mi}{2}\PY{+w}{  }\PY{l+m+mi}{4}\PY{+w}{  }\PY{l+m+mi}{2}\PY{+w}{   }\PY{l+m+mi}{2}\PY{+w}{  }\PY{l+m+mi}{4}\PY{p}{;}\PY{+w}{  }\PY{c+cm}{\PYZsh{}=}
\PY{c+cm}{ }\PY{c+cm}{ }\PY{c+cm}{ }\PY{c+cm}{ }\PY{c+cm}{ }\PY{c+cm}{ }\PY{c+cm}{=\PYZsh{}}\PY{+w}{ }\PY{l+m+mi}{9}\PY{+w}{  }\PY{l+m+mi}{5}\PY{+w}{  }\PY{l+m+mi}{9}\PY{+w}{   }\PY{l+m+mi}{8}\PY{+w}{  }\PY{l+m+mi}{7}\PY{p}{;}\PY{+w}{  }\PY{c+cm}{\PYZsh{}=}
\PY{c+cm}{ }\PY{c+cm}{ }\PY{c+cm}{ }\PY{c+cm}{ }\PY{c+cm}{ }\PY{c+cm}{ }\PY{c+cm}{=\PYZsh{}}\PY{+w}{ }\PY{l+m+mi}{9}\PY{+w}{  }\PY{l+m+mi}{3}\PY{+w}{  }\PY{l+m+mi}{10}\PY{+w}{  }\PY{l+m+mi}{2}\PY{+w}{  }\PY{l+m+mi}{2}\PY{p}{;}\PY{+w}{  }\PY{c+cm}{\PYZsh{}=}
\PY{c+cm}{ }\PY{c+cm}{ }\PY{c+cm}{ }\PY{c+cm}{ }\PY{c+cm}{ }\PY{c+cm}{ }\PY{c+cm}{=\PYZsh{}}\PY{+w}{ }\PY{l+m+mi}{1}\PY{+w}{  }\PY{l+m+mi}{2}\PY{+w}{  }\PY{l+m+mi}{10}\PY{+w}{  }\PY{l+m+mi}{7}\PY{+w}{  }\PY{l+m+mi}{8}\PY{p}{;}\PY{+w}{  }\PY{c+cm}{\PYZsh{}=}
\PY{c+cm}{ }\PY{c+cm}{ }\PY{c+cm}{ }\PY{c+cm}{ }\PY{c+cm}{ }\PY{c+cm}{ }\PY{c+cm}{=\PYZsh{}}\PY{+w}{ }\PY{l+m+mi}{4}\PY{+w}{  }\PY{l+m+mi}{3}\PY{+w}{  }\PY{l+m+mi}{3}\PY{+w}{   }\PY{l+m+mi}{8}\PY{+w}{  }\PY{l+m+mi}{6}\PY{+w}{   }\PY{c}{\PYZsh{}}
\PY{+w}{   }\PY{p}{]}\PY{p}{;}

\PY{+w}{   }\PY{n}{A2}\PY{+w}{ }\PY{o}{=}\PY{+w}{ }\PY{p}{[}\PY{+w}{                 }\PY{c+cm}{\PYZsh{}=}
\PY{c+cm}{ }\PY{c+cm}{ }\PY{c+cm}{ }\PY{c+cm}{ }\PY{c+cm}{ }\PY{c+cm}{ }\PY{c+cm}{=\PYZsh{}}\PY{+w}{  }\PY{l+m+mi}{1}\PY{+w}{  }\PY{o}{\PYZhy{}}\PY{l+m+mi}{2}\PY{+w}{  }\PY{o}{\PYZhy{}}\PY{l+m+mi}{2}\PY{+w}{  }\PY{l+m+mi}{1}\PY{+w}{ }\PY{o}{\PYZhy{}}\PY{l+m+mi}{4}\PY{p}{;}\PY{+w}{  }\PY{c+cm}{\PYZsh{}=}
\PY{c+cm}{ }\PY{c+cm}{ }\PY{c+cm}{ }\PY{c+cm}{ }\PY{c+cm}{ }\PY{c+cm}{ }\PY{c+cm}{=\PYZsh{}}\PY{+w}{ }\PY{o}{\PYZhy{}}\PY{l+m+mi}{1}\PY{+w}{  }\PY{o}{\PYZhy{}}\PY{l+m+mi}{3}\PY{+w}{   }\PY{l+m+mi}{1}\PY{+w}{  }\PY{l+m+mi}{5}\PY{+w}{  }\PY{l+m+mi}{3}\PY{p}{;}\PY{+w}{  }\PY{c+cm}{\PYZsh{}=}
\PY{c+cm}{ }\PY{c+cm}{ }\PY{c+cm}{ }\PY{c+cm}{ }\PY{c+cm}{ }\PY{c+cm}{ }\PY{c+cm}{=\PYZsh{}}\PY{+w}{  }\PY{l+m+mi}{3}\PY{+w}{  }\PY{o}{\PYZhy{}}\PY{l+m+mi}{2}\PY{+w}{   }\PY{l+m+mi}{0}\PY{+w}{  }\PY{l+m+mi}{5}\PY{+w}{ }\PY{o}{\PYZhy{}}\PY{l+m+mi}{3}\PY{p}{;}\PY{+w}{  }\PY{c+cm}{\PYZsh{}=}
\PY{c+cm}{ }\PY{c+cm}{ }\PY{c+cm}{ }\PY{c+cm}{ }\PY{c+cm}{ }\PY{c+cm}{ }\PY{c+cm}{=\PYZsh{}}\PY{+w}{  }\PY{l+m+mi}{0}\PY{+w}{  }\PY{o}{\PYZhy{}}\PY{l+m+mi}{1}\PY{+w}{  }\PY{o}{\PYZhy{}}\PY{l+m+mi}{2}\PY{+w}{ }\PY{o}{\PYZhy{}}\PY{l+m+mi}{3}\PY{+w}{ }\PY{o}{\PYZhy{}}\PY{l+m+mi}{3}\PY{p}{;}\PY{+w}{  }\PY{c+cm}{\PYZsh{}=}
\PY{c+cm}{ }\PY{c+cm}{ }\PY{c+cm}{ }\PY{c+cm}{ }\PY{c+cm}{ }\PY{c+cm}{ }\PY{c+cm}{=\PYZsh{}}\PY{+w}{ }\PY{o}{\PYZhy{}}\PY{l+m+mi}{2}\PY{+w}{  }\PY{o}{\PYZhy{}}\PY{l+m+mi}{1}\PY{+w}{  }\PY{o}{\PYZhy{}}\PY{l+m+mi}{4}\PY{+w}{ }\PY{o}{\PYZhy{}}\PY{l+m+mi}{5}\PY{+w}{ }\PY{o}{\PYZhy{}}\PY{l+m+mi}{4}\PY{+w}{   }\PY{c}{\PYZsh{}}
\PY{+w}{   }\PY{p}{]}\PY{p}{;}

\PY{+w}{   }\PY{n}{println}\PY{p}{(}\PY{l+s}{\PYZdq{}}\PY{l+s}{A1: det: }\PY{l+s}{\PYZdq{}}\PY{p}{,}\PY{+w}{ }\PY{n}{det}\PY{p}{(}\PY{n}{A1}\PY{p}{)}\PY{p}{)}
\PY{+w}{   }\PY{n}{println}\PY{p}{(}\PY{l+s}{\PYZdq{}}\PY{l+s}{    our result: }\PY{l+s}{\PYZdq{}}\PY{p}{,}\PY{+w}{ }\PY{n}{determinant}\PY{p}{(}\PY{n}{A1}\PY{p}{)}\PY{p}{)}
\PY{+w}{   }\PY{n}{println}\PY{p}{(}\PY{l+s}{\PYZdq{}}\PY{l+s}{    error: }\PY{l+s}{\PYZdq{}}\PY{p}{,}\PY{+w}{ }\PY{n}{abs}\PY{p}{(}\PY{n}{determinant}\PY{p}{(}\PY{n}{A1}\PY{p}{)}\PY{+w}{ }\PY{o}{\PYZhy{}}\PY{+w}{ }\PY{n}{det}\PY{p}{(}\PY{n}{A1}\PY{p}{)}\PY{p}{)}\PY{p}{)}

\PY{+w}{   }\PY{n}{println}\PY{p}{(}\PY{l+s}{\PYZdq{}}\PY{l+s}{A2: det: }\PY{l+s}{\PYZdq{}}\PY{p}{,}\PY{+w}{ }\PY{n}{det}\PY{p}{(}\PY{n}{A2}\PY{p}{)}\PY{p}{)}
\PY{+w}{   }\PY{n}{println}\PY{p}{(}\PY{l+s}{\PYZdq{}}\PY{l+s}{    our result: }\PY{l+s}{\PYZdq{}}\PY{p}{,}\PY{+w}{ }\PY{n}{determinant}\PY{p}{(}\PY{n}{A2}\PY{p}{)}\PY{p}{)}
\PY{+w}{   }\PY{n}{println}\PY{p}{(}\PY{l+s}{\PYZdq{}}\PY{l+s}{    error: }\PY{l+s}{\PYZdq{}}\PY{p}{,}\PY{+w}{ }\PY{n}{abs}\PY{p}{(}\PY{n}{determinant}\PY{p}{(}\PY{n}{A2}\PY{p}{)}\PY{+w}{ }\PY{o}{\PYZhy{}}\PY{+w}{ }\PY{n}{det}\PY{p}{(}\PY{n}{A2}\PY{p}{)}\PY{p}{)}\PY{p}{)}

\PY{k}{end}
\end{Verbatim}
\end{tcolorbox}

    \begin{Verbatim}[commandchars=\\\{\}]
A1: det: 281.9999999999995
    our result: 282.0000000000009
    error: 1.4210854715202004e-12
A2: det: 2.000000000000068
    our result: 1.999999999999907
    error: 1.609823385706477e-13
    \end{Verbatim}

    As we can see, the absolute error (also relative error since their
determinants \(> 1\)) is extremely small.

    \hypertarget{section-2.5-8}{%
\subsection{Section 2.5 \#8}\label{section-2.5-8}}

    The UpperTriangular and LowerTriangular matrix types cause specialized
algorithms to be invoked by the backslash. Define

    \begin{tcolorbox}[breakable, size=fbox, boxrule=1pt, pad at break*=1mm,colback=cellbackground, colframe=cellborder]
\prompt{In}{incolor}{217}{\boxspacing}
\begin{Verbatim}[commandchars=\\\{\}]
\PY{k}{using}\PY{+w}{ }\PY{n}{LinearAlgebra}

\PY{n}{A}\PY{+w}{ }\PY{o}{=}\PY{+w}{ }\PY{n}{rand}\PY{p}{(}\PY{l+m+mi}{1000}\PY{p}{,}\PY{+w}{ }\PY{l+m+mi}{1000}\PY{p}{)}
\PY{n}{B}\PY{+w}{ }\PY{o}{=}\PY{+w}{ }\PY{n}{tril}\PY{p}{(}\PY{n}{A}\PY{p}{)}
\PY{n}{C}\PY{+w}{ }\PY{o}{=}\PY{+w}{ }\PY{n}{LowerTriangular}\PY{p}{(}\PY{n}{B}\PY{p}{)}
\PY{n}{b}\PY{+w}{ }\PY{o}{=}\PY{+w}{ }\PY{n}{rand}\PY{p}{(}\PY{l+m+mi}{1000}\PY{p}{)}\PY{p}{;}
\end{Verbatim}
\end{tcolorbox}

    Using \texttt{@elapsed} with the backslash solver, time how long it
takes to solve the linear system \(Ax=b\) 100 times; then, do the same
for matrices \(B\) and \(C\). Is the timing for \(B\) closer to \(A\) or
to \(C\)? (Hint: Remember to make one timing run without recording
results, so that compilation time is not counted.)

As a bonus, check the time if you use the LU factored form of A.

Hint: this determines the length of time to compute sine of 10 million
Normally-distributed random numbers. (About 0.4 s on my computer.) The
\texttt{\_} is a digits separator used to help humans read numbers with
many digits.

    \begin{tcolorbox}[breakable, size=fbox, boxrule=1pt, pad at break*=1mm,colback=cellbackground, colframe=cellborder]
\prompt{In}{incolor}{218}{\boxspacing}
\begin{Verbatim}[commandchars=\\\{\}]
\PY{n}{println}\PY{p}{(}\PY{l+s}{\PYZdq{}}\PY{l+s}{Using backslash: }\PY{l+s}{\PYZdq{}}\PY{p}{)}
\PY{n}{println}\PY{p}{(}\PY{l+s}{\PYZdq{}}\PY{l+s}{ A: }\PY{l+s}{\PYZdq{}}\PY{p}{,}\PY{+w}{ }\PY{n+nd}{@elapsed}\PY{+w}{ }\PY{p}{(}\PY{k}{for}\PY{+w}{ }\PY{n}{i}\PY{+w}{ }\PY{k}{in}\PY{+w}{ }\PY{l+m+mi}{1}\PY{+w}{ }\PY{o}{:}\PY{+w}{ }\PY{l+m+mi}{100}\PY{p}{;}\PY{+w}{ }
\PY{+w}{    }\PY{n}{A}\PY{+w}{ }\PY{o}{\PYZbs{}}\PY{+w}{ }\PY{n}{b}\PY{p}{;}
\PY{k}{end}\PY{p}{)}\PY{p}{)}
\PY{n}{println}\PY{p}{(}\PY{l+s}{\PYZdq{}}\PY{l+s}{ B: }\PY{l+s}{\PYZdq{}}\PY{p}{,}\PY{+w}{ }\PY{n+nd}{@elapsed}\PY{+w}{ }\PY{p}{(}\PY{k}{for}\PY{+w}{ }\PY{n}{i}\PY{+w}{ }\PY{k}{in}\PY{+w}{ }\PY{l+m+mi}{1}\PY{+w}{ }\PY{o}{:}\PY{+w}{ }\PY{l+m+mi}{100}\PY{p}{;}\PY{+w}{ }
\PY{+w}{    }\PY{n}{B}\PY{+w}{ }\PY{o}{\PYZbs{}}\PY{+w}{ }\PY{n}{b}\PY{p}{;}
\PY{k}{end}\PY{p}{)}\PY{p}{)}
\PY{n}{println}\PY{p}{(}\PY{l+s}{\PYZdq{}}\PY{l+s}{ C: }\PY{l+s}{\PYZdq{}}\PY{p}{,}\PY{+w}{ }\PY{n+nd}{@elapsed}\PY{+w}{ }\PY{p}{(}\PY{k}{for}\PY{+w}{ }\PY{n}{i}\PY{+w}{ }\PY{k}{in}\PY{+w}{ }\PY{l+m+mi}{1}\PY{+w}{ }\PY{o}{:}\PY{+w}{ }\PY{l+m+mi}{100}\PY{p}{;}\PY{+w}{ }
\PY{+w}{    }\PY{n}{C}\PY{+w}{ }\PY{o}{\PYZbs{}}\PY{+w}{ }\PY{n}{b}\PY{p}{;}
\PY{k}{end}\PY{p}{)}\PY{p}{)}
\end{Verbatim}
\end{tcolorbox}

    \begin{Verbatim}[commandchars=\\\{\}]
Using backslash:
 A: 2.2112757
 B: 0.041298
 C: 0.0092961
    \end{Verbatim}

    The timing for \(B\) is closer to \(C\) rather than \(B\), this is
probably because that \(B, C\) are all triangular matrices, and there
exist some specified methods for such matrices.

    \begin{tcolorbox}[breakable, size=fbox, boxrule=1pt, pad at break*=1mm,colback=cellbackground, colframe=cellborder]
\prompt{In}{incolor}{219}{\boxspacing}
\begin{Verbatim}[commandchars=\\\{\}]
\PY{n}{println}\PY{p}{(}\PY{l+s}{\PYZdq{}}\PY{l+s}{ A with LU decomposition: }\PY{l+s}{\PYZdq{}}\PY{p}{,}\PY{+w}{ }\PY{n+nd}{@elapsed}\PY{+w}{ }\PY{p}{(}\PY{k}{for}\PY{+w}{ }\PY{n}{i}\PY{+w}{ }\PY{k}{in}\PY{+w}{ }\PY{l+m+mi}{1}\PY{+w}{ }\PY{o}{:}\PY{+w}{ }\PY{l+m+mi}{100}\PY{p}{;}\PY{+w}{ }
\PY{+w}{    }\PY{n}{L}\PY{p}{,}\PY{+w}{ }\PY{n}{U}\PY{+w}{ }\PY{o}{=}\PY{+w}{ }\PY{n}{lu}\PY{p}{(}\PY{n}{A}\PY{p}{)}\PY{p}{;}
\PY{+w}{    }\PY{n}{y}\PY{+w}{ }\PY{o}{=}\PY{+w}{ }\PY{n}{L}\PY{+w}{ }\PY{o}{\PYZbs{}}\PY{+w}{ }\PY{n}{b}\PY{p}{;}
\PY{+w}{    }\PY{n}{U}\PY{+w}{ }\PY{o}{\PYZbs{}}\PY{+w}{ }\PY{n}{y}
\PY{k}{end}\PY{p}{)}\PY{p}{)}

\PY{n}{L}\PY{+w}{ }\PY{o}{=}\PY{+w}{ }\PY{n+nb}{nothing}\PY{p}{;}\PY{+w}{ }\PY{n}{U}\PY{+w}{ }\PY{o}{=}\PY{+w}{ }\PY{n+nb}{nothing}\PY{p}{;}
\PY{n}{println}\PY{p}{(}\PY{l+s}{\PYZdq{}}\PY{l+s}{ A with LU pre\PYZhy{}decomposition: }\PY{l+s}{\PYZdq{}}\PY{p}{,}\PY{+w}{ }\PY{n+nd}{@elapsed}\PY{+w}{ }\PY{p}{(}\PY{k}{for}\PY{+w}{ }\PY{n}{i}\PY{+w}{ }\PY{k}{in}\PY{+w}{ }\PY{l+m+mi}{1}\PY{+w}{ }\PY{o}{:}\PY{+w}{ }\PY{l+m+mi}{100}\PY{p}{;}\PY{+w}{ }
\PY{+w}{    }\PY{k}{if}\PY{+w}{ }\PY{p}{(}\PY{n}{i}\PY{+w}{ }\PY{o}{==}\PY{+w}{ }\PY{l+m+mi}{1}\PY{p}{)}
\PY{+w}{        }\PY{n}{L}\PY{p}{,}\PY{+w}{ }\PY{n}{U}\PY{+w}{ }\PY{o}{=}\PY{+w}{ }\PY{n}{lu}\PY{p}{(}\PY{n}{A}\PY{p}{)}\PY{p}{;}
\PY{+w}{    }\PY{k}{end}
\PY{+w}{    }\PY{n}{y}\PY{+w}{ }\PY{o}{=}\PY{+w}{ }\PY{n}{L}\PY{+w}{ }\PY{o}{\PYZbs{}}\PY{+w}{ }\PY{n}{b}\PY{p}{;}
\PY{+w}{    }\PY{n}{U}\PY{+w}{ }\PY{o}{\PYZbs{}}\PY{+w}{ }\PY{n}{y}
\PY{k}{end}\PY{p}{)}\PY{p}{)}
\end{Verbatim}
\end{tcolorbox}

    \begin{Verbatim}[commandchars=\\\{\}]
 A with LU decomposition: 3.2957766
 A with LU pre-decomposition: 0.2132572
    \end{Verbatim}

    Solving time increased by \(0.4\) s if we decompose \(A\) to \(L, U\)
every time then solve.

But if we decompose it once, then solve the equation for \(100\) times
reusing \(L, U\), the time decreases hugely to \(0.13\) s.

    \hypertarget{section-2.9-5}{%
\subsection{Section 2.9 \#5}\label{section-2.9-5}}

The Tridiagonal matrix type invokes a specialized algorithm for solving
a linear system.

\begin{enumerate}
\def\labelenumi{(\alph{enumi})}
\tightlist
\item
  Set n=1000 and t=0. In a loop that runs 50 times, generate a linear
  system via
\end{enumerate}

    \begin{tcolorbox}[breakable, size=fbox, boxrule=1pt, pad at break*=1mm,colback=cellbackground, colframe=cellborder]
\prompt{In}{incolor}{220}{\boxspacing}
\begin{Verbatim}[commandchars=\\\{\}]
\PY{n}{n1}\PY{+w}{ }\PY{o}{=}\PY{+w}{ }\PY{l+m+mi}{1000}
\PY{n}{A1}\PY{+w}{ }\PY{o}{=}\PY{+w}{ }\PY{n}{triu}\PY{p}{(}\PY{+w}{ }\PY{n}{tril}\PY{p}{(}\PY{n}{rand}\PY{p}{(}\PY{n}{n}\PY{p}{,}\PY{n}{n}\PY{p}{)}\PY{p}{,}\PY{l+m+mi}{1}\PY{p}{)}\PY{p}{,}\PY{+w}{ }\PY{o}{\PYZhy{}}\PY{l+m+mi}{1}\PY{p}{)}
\PY{n}{b1}\PY{+w}{ }\PY{o}{=}\PY{+w}{ }\PY{n}{ones}\PY{p}{(}\PY{n}{n}\PY{p}{)}\PY{p}{;}

\PY{n}{ts1}\PY{+w}{ }\PY{o}{=}\PY{+w}{ }\PY{p}{[}\PY{p}{(}\PY{n+nd}{@elapsed}\PY{+w}{ }\PY{p}{(}\PY{k}{for}\PY{+w}{ }\PY{n}{i}\PY{+w}{ }\PY{k}{in}\PY{+w}{ }\PY{l+m+mi}{1}\PY{+w}{ }\PY{o}{:}\PY{+w}{ }\PY{n}{t}\PY{p}{;}\PY{+w}{ }
\PY{n}{A1}\PY{+w}{ }\PY{o}{\PYZbs{}}\PY{+w}{ }\PY{n}{b1}\PY{p}{;}
\PY{k}{end}\PY{p}{)}\PY{p}{)}\PY{+w}{ }\PY{k}{for}\PY{+w}{ }\PY{n}{t}\PY{+w}{ }\PY{k}{in}\PY{+w}{ }\PY{n}{range}\PY{p}{(}\PY{l+m+mi}{0}\PY{p}{,}\PY{+w}{ }\PY{l+m+mi}{50}\PY{p}{)}\PY{p}{]}
\PY{n}{println}\PY{p}{(}\PY{l+s}{\PYZdq{}}\PY{l+s}{Final value of t: }\PY{l+s}{\PYZdq{}}\PY{p}{,}\PY{+w}{ }\PY{n}{ts1}\PY{p}{[}\PY{k}{end}\PY{p}{]}\PY{p}{)}
\end{Verbatim}
\end{tcolorbox}

    \begin{Verbatim}[commandchars=\\\{\}]
Final value of t: 1.1899208
    \end{Verbatim}

    \begin{tcolorbox}[breakable, size=fbox, boxrule=1pt, pad at break*=1mm,colback=cellbackground, colframe=cellborder]
\prompt{In}{incolor}{221}{\boxspacing}
\begin{Verbatim}[commandchars=\\\{\}]
\PY{n}{plot}\PY{p}{(}\PY{n}{ts1}\PY{p}{,}\PY{+w}{ }\PY{n}{label}\PY{o}{=}\PY{l+s}{\PYZdq{}}\PY{l+s}{Elapsed Time}\PY{l+s}{\PYZdq{}}\PY{p}{,}\PY{+w}{ }\PY{n}{xlabel}\PY{o}{=}\PY{l+s}{\PYZdq{}}\PY{l+s}{Count}\PY{l+s}{\PYZdq{}}\PY{p}{,}\PY{+w}{ }\PY{n}{ylabel}\PY{o}{=}\PY{l+s}{\PYZdq{}}\PY{l+s}{Time (s)}\PY{l+s}{\PYZdq{}}\PY{p}{)}
\end{Verbatim}
\end{tcolorbox}

    \begin{center}
    \adjustimage{max size={0.9\linewidth}{0.9\paperheight}}{hw-02_files/hw-02_35_0.pdf}
    \end{center}
    { \hspace*{\fill} \\}
    
    Using @elapsed, increment t by the time it takes to perform A\b. Print
out the final value of t.

\begin{enumerate}
\def\labelenumi{(\alph{enumi})}
\setcounter{enumi}{1}
\tightlist
\item
  Repeat the experiment of part (a), but generate the matrix via
  \texttt{A\ =\ Tridiagonal(rand(n,n))} What is the ratio of running
  times for part (a) and (b)?
\end{enumerate}

    \begin{tcolorbox}[breakable, size=fbox, boxrule=1pt, pad at break*=1mm,colback=cellbackground, colframe=cellborder]
\prompt{In}{incolor}{222}{\boxspacing}
\begin{Verbatim}[commandchars=\\\{\}]
\PY{n}{n2}\PY{+w}{ }\PY{o}{=}\PY{+w}{ }\PY{l+m+mi}{1000}
\PY{n}{A2}\PY{+w}{ }\PY{o}{=}\PY{+w}{ }\PY{n}{Tridiagonal}\PY{p}{(}\PY{n}{rand}\PY{p}{(}\PY{n}{n2}\PY{p}{,}\PY{+w}{ }\PY{n}{n2}\PY{p}{)}\PY{p}{)}
\PY{n}{b2}\PY{+w}{ }\PY{o}{=}\PY{+w}{ }\PY{n}{ones}\PY{p}{(}\PY{n}{n}\PY{p}{)}\PY{p}{;}

\PY{n}{ts2}\PY{+w}{ }\PY{o}{=}\PY{+w}{ }\PY{p}{[}\PY{p}{(}\PY{n+nd}{@elapsed}\PY{+w}{ }\PY{p}{(}\PY{k}{for}\PY{+w}{ }\PY{n}{i}\PY{+w}{ }\PY{k}{in}\PY{+w}{ }\PY{l+m+mi}{1}\PY{+w}{ }\PY{o}{:}\PY{+w}{ }\PY{n}{t}\PY{p}{;}\PY{+w}{ }
\PY{n}{A2}\PY{+w}{ }\PY{o}{\PYZbs{}}\PY{+w}{ }\PY{n}{b2}\PY{p}{;}
\PY{k}{end}\PY{p}{)}\PY{p}{)}\PY{+w}{ }\PY{k}{for}\PY{+w}{ }\PY{n}{t}\PY{+w}{ }\PY{k}{in}\PY{+w}{ }\PY{n}{range}\PY{p}{(}\PY{l+m+mi}{0}\PY{p}{,}\PY{+w}{ }\PY{l+m+mi}{50}\PY{p}{)}\PY{p}{]}
\PY{n}{println}\PY{p}{(}\PY{l+s}{\PYZdq{}}\PY{l+s}{Final value of t: }\PY{l+s}{\PYZdq{}}\PY{p}{,}\PY{+w}{ }\PY{n}{ts2}\PY{p}{[}\PY{k}{end}\PY{p}{]}\PY{p}{)}
\end{Verbatim}
\end{tcolorbox}

    \begin{Verbatim}[commandchars=\\\{\}]
Final value of t: 0.0013423
    \end{Verbatim}

    \begin{tcolorbox}[breakable, size=fbox, boxrule=1pt, pad at break*=1mm,colback=cellbackground, colframe=cellborder]
\prompt{In}{incolor}{223}{\boxspacing}
\begin{Verbatim}[commandchars=\\\{\}]
\PY{n}{plot}\PY{p}{(}\PY{n}{ts2}\PY{p}{,}\PY{+w}{ }\PY{n}{label}\PY{o}{=}\PY{l+s}{\PYZdq{}}\PY{l+s}{Elapsed Time}\PY{l+s}{\PYZdq{}}\PY{p}{,}\PY{+w}{ }\PY{n}{xlabel}\PY{o}{=}\PY{l+s}{\PYZdq{}}\PY{l+s}{Count}\PY{l+s}{\PYZdq{}}\PY{p}{,}\PY{+w}{ }\PY{n}{ylabel}\PY{o}{=}\PY{l+s}{\PYZdq{}}\PY{l+s}{Time (s)}\PY{l+s}{\PYZdq{}}\PY{p}{,}\PY{+w}{ }\PY{n}{title}\PY{o}{=}\PY{l+s}{\PYZdq{}}\PY{l+s}{(Triangular)}\PY{l+s}{\PYZdq{}}\PY{p}{)}
\end{Verbatim}
\end{tcolorbox}

    \begin{center}
    \adjustimage{max size={0.9\linewidth}{0.9\paperheight}}{hw-02_files/hw-02_38_0.pdf}
    \end{center}
    { \hspace*{\fill} \\}
    
    \begin{tcolorbox}[breakable, size=fbox, boxrule=1pt, pad at break*=1mm,colback=cellbackground, colframe=cellborder]
\prompt{In}{incolor}{224}{\boxspacing}
\begin{Verbatim}[commandchars=\\\{\}]
\PY{n}{plot}\PY{p}{(}\PY{n}{ts1}\PY{+w}{ }\PY{o}{./}\PY{+w}{ }\PY{n}{ts2}\PY{p}{,}\PY{+w}{ }\PY{n}{label}\PY{o}{=}\PY{n+nb}{nothing}\PY{p}{,}\PY{+w}{ }\PY{n}{xlabel}\PY{o}{=}\PY{l+s}{\PYZdq{}}\PY{l+s}{Count}\PY{l+s}{\PYZdq{}}\PY{p}{,}\PY{+w}{ }\PY{n}{ylabel}\PY{o}{=}\PY{l+s}{\PYZdq{}}\PY{l+s}{Ratio}\PY{l+s}{\PYZdq{}}\PY{p}{,}\PY{+w}{ }\PY{n}{title}\PY{o}{=}\PY{l+s}{\PYZdq{}}\PY{l+s}{Time ratio}\PY{l+s}{\PYZdq{}}\PY{p}{)}
\end{Verbatim}
\end{tcolorbox}

    \begin{center}
    \adjustimage{max size={0.9\linewidth}{0.9\paperheight}}{hw-02_files/hw-02_39_0.pdf}
    \end{center}
    { \hspace*{\fill} \\}
    
    \begin{tcolorbox}[breakable, size=fbox, boxrule=1pt, pad at break*=1mm,colback=cellbackground, colframe=cellborder]
\prompt{In}{incolor}{225}{\boxspacing}
\begin{Verbatim}[commandchars=\\\{\}]
\PY{n}{sum}\PY{p}{(}\PY{p}{(}\PY{n}{ts1}\PY{+w}{ }\PY{o}{./}\PY{+w}{ }\PY{n}{ts2}\PY{p}{)}\PY{p}{[}\PY{l+m+mi}{5}\PY{o}{:}\PY{k}{end}\PY{p}{]}\PY{p}{)}\PY{+w}{ }\PY{o}{/}\PY{+w}{ }\PY{p}{(}\PY{n}{length}\PY{p}{(}\PY{n}{ts1}\PY{p}{)}\PY{+w}{ }\PY{o}{\PYZhy{}}\PY{+w}{ }\PY{l+m+mi}{5}\PY{+w}{ }\PY{o}{+}\PY{+w}{ }\PY{l+m+mi}{1}\PY{p}{)}
\end{Verbatim}
\end{tcolorbox}

    
    \begin{Verbatim}[commandchars=\\\{\}]
778.8295739569452
    \end{Verbatim}

    
    From the graph, we can notice that from count=5, the ratio became
relatively stable, so we discarded indices 0\textasciitilde4, and
calculated the average ratio, it's about \(646\).

    \begin{enumerate}
\def\labelenumi{(\alph{enumi})}
\setcounter{enumi}{2}
\tightlist
\item
  Now perform the experiment of part (b) for \(n = 1000 + 200k\),
  \(k = 0, \dots, 15\), keeping the total time for each value of \(n\)
  in a vector. Plot running time as a function of \(n\) on a log-log
  scale. Is the time most like \(O(n)\), \(O(n^2)\), or \(O(n^3)\)? (If
  the answer is unclear, try increasing the number of solves per value
  of to 100 or more.)
\end{enumerate}

    \begin{tcolorbox}[breakable, size=fbox, boxrule=1pt, pad at break*=1mm,colback=cellbackground, colframe=cellborder]
\prompt{In}{incolor}{226}{\boxspacing}
\begin{Verbatim}[commandchars=\\\{\}]
\PY{n}{ns}\PY{+w}{ }\PY{o}{=}\PY{+w}{ }\PY{l+m+mi}{1000}\PY{+w}{ }\PY{o}{.+}\PY{+w}{ }\PY{l+m+mi}{200}\PY{+w}{ }\PY{o}{.*}\PY{+w}{ }\PY{p}{(}\PY{l+m+mi}{0}\PY{o}{:}\PY{l+m+mi}{15}\PY{p}{)}

\PY{n}{ts3}\PY{+w}{ }\PY{o}{=}\PY{+w}{ }\PY{p}{[}\PY{l+m+mi}{0}\PY{p}{]}
\PY{k}{for}\PY{+w}{ }\PY{n}{n}\PY{+w}{ }\PY{k}{in}\PY{+w}{ }\PY{n}{ns}\PY{p}{;}
\PY{+w}{    }\PY{n}{C}\PY{+w}{ }\PY{o}{=}\PY{+w}{ }\PY{n}{Tridiagonal}\PY{p}{(}\PY{n}{rand}\PY{p}{(}\PY{n}{n}\PY{p}{,}\PY{+w}{ }\PY{n}{n}\PY{p}{)}\PY{p}{)}
\PY{+w}{    }\PY{n}{b3}\PY{+w}{ }\PY{o}{=}\PY{+w}{ }\PY{n}{ones}\PY{p}{(}\PY{n}{n}\PY{p}{)}
\PY{+w}{    }\PY{n}{timi}\PY{+w}{ }\PY{o}{=}\PY{+w}{ }\PY{p}{(}\PY{n+nd}{@elapsed}\PY{+w}{ }\PY{p}{(}\PY{k}{for}\PY{+w}{ }\PY{n}{i}\PY{+w}{ }\PY{k}{in}\PY{+w}{ }\PY{l+m+mi}{1}\PY{+w}{ }\PY{o}{:}\PY{+w}{ }\PY{l+m+mi}{10000}\PY{p}{;}\PY{+w}{ }
\PY{+w}{                }\PY{n}{C}\PY{+w}{ }\PY{o}{\PYZbs{}}\PY{+w}{ }\PY{n}{b3}\PY{p}{;}
\PY{+w}{            }\PY{k}{end}\PY{p}{)}\PY{p}{)}
\PY{+w}{    }\PY{n}{ts3}\PY{+w}{ }\PY{o}{=}\PY{+w}{ }\PY{p}{[}\PY{n}{ts3}\PY{+w}{ }\PY{n}{timi}\PY{p}{]}\PY{p}{;}
\PY{k}{end}
\PY{n}{ts3}\PY{+w}{ }\PY{o}{=}\PY{+w}{ }\PY{n}{ts3}\PY{p}{[}\PY{l+m+mi}{2}\PY{o}{:}\PY{k}{end}\PY{p}{]}\PY{p}{;}
\end{Verbatim}
\end{tcolorbox}

    \begin{tcolorbox}[breakable, size=fbox, boxrule=1pt, pad at break*=1mm,colback=cellbackground, colframe=cellborder]
\prompt{In}{incolor}{227}{\boxspacing}
\begin{Verbatim}[commandchars=\\\{\}]
\PY{n}{plot}\PY{p}{(}\PY{n}{ns}\PY{p}{,}\PY{+w}{ }\PY{n}{ts3}\PY{p}{,}\PY{+w}{ }\PY{n}{xaxis}\PY{o}{=}\PY{l+s+ss}{:log}\PY{p}{,}\PY{+w}{ }\PY{n}{yaxis}\PY{o}{=}\PY{l+s+ss}{:log}\PY{p}{,}
\PY{+w}{    }\PY{n}{xlabel}\PY{o}{=}\PY{l+s}{\PYZdq{}}\PY{l+s}{n}\PY{l+s}{\PYZdq{}}\PY{p}{,}\PY{+w}{ }\PY{n}{ylabel}\PY{o}{=}\PY{l+s}{\PYZdq{}}\PY{l+s}{Time}\PY{l+s}{\PYZdq{}}\PY{p}{,}\PY{+w}{ }\PY{n}{label}\PY{o}{=}\PY{n+nb}{nothing}
\PY{p}{)}
\end{Verbatim}
\end{tcolorbox}

    \begin{center}
    \adjustimage{max size={0.9\linewidth}{0.9\paperheight}}{hw-02_files/hw-02_44_0.pdf}
    \end{center}
    { \hspace*{\fill} \\}
    
    The graph is nearly a straight line, hence the time complexity is most
likely to be \(O(n)\).

    \hypertarget{section-3.1-2}{%
\subsection{Section 3.1 \#2}\label{section-3.1-2}}

Here are counts of the U.S. population in millions from the census
performed every ten years, beginning with 1790 and ending with 2010.

    \begin{tcolorbox}[breakable, size=fbox, boxrule=1pt, pad at break*=1mm,colback=cellbackground, colframe=cellborder]
\prompt{In}{incolor}{228}{\boxspacing}
\begin{Verbatim}[commandchars=\\\{\}]
\PY{n}{usa\PYZus{}pop}\PY{+w}{ }\PY{o}{=}\PY{+w}{ }\PY{p}{[}\PY{l+m+mf}{3.929}\PY{p}{,}\PY{+w}{ }\PY{l+m+mf}{5.308}\PY{p}{,}\PY{+w}{ }\PY{l+m+mf}{7.240}\PY{p}{,}\PY{+w}{ }\PY{l+m+mf}{9.638}\PY{p}{,}\PY{+w}{ }\PY{l+m+mf}{12.87}\PY{p}{,}\PY{+w}{ }\PY{l+m+mf}{17.07}\PY{p}{,}\PY{+w}{ }\PY{l+m+mf}{23.19}\PY{p}{,}\PY{+w}{ }\PY{l+m+mf}{31.44}\PY{p}{,}\PY{+w}{ }\PY{l+m+mf}{39.82}\PY{p}{,}\PY{+w}{ }\PY{l+m+mf}{50.19}\PY{p}{,}\PY{+w}{ }\PY{l+m+mf}{62.95}\PY{p}{,}\PY{+w}{ }\PY{l+m+mf}{76.21}\PY{p}{,}
\PY{l+m+mf}{92.22}\PY{p}{,}\PY{+w}{ }\PY{l+m+mf}{106.0}\PY{p}{,}\PY{+w}{ }\PY{l+m+mf}{122.8}\PY{p}{,}\PY{+w}{ }\PY{l+m+mf}{132.2}\PY{p}{,}\PY{+w}{ }\PY{l+m+mf}{150.7}\PY{p}{,}\PY{+w}{ }\PY{l+m+mf}{179.3}\PY{p}{,}\PY{+w}{ }\PY{l+m+mf}{203.3}\PY{p}{,}\PY{+w}{ }\PY{l+m+mf}{226.5}\PY{p}{,}\PY{+w}{ }\PY{l+m+mf}{248.7}\PY{p}{,}\PY{+w}{ }\PY{l+m+mf}{281.4}\PY{p}{,}\PY{+w}{ }\PY{l+m+mf}{308.7}\PY{p}{]}\PY{p}{;}
\end{Verbatim}
\end{tcolorbox}

    \begin{enumerate}
\def\labelenumi{(\alph{enumi})}
\tightlist
\item
  Find a best-fitting cubic polynomial for the data. Plot the data as
  points superimposed on a (smooth) graph of the cubic over the full
  range of time. Label the axes. What does the fit predict for the
  population in the years 2000, 2010, and 2020?
\end{enumerate}

    \begin{tcolorbox}[breakable, size=fbox, boxrule=1pt, pad at break*=1mm,colback=cellbackground, colframe=cellborder]
\prompt{In}{incolor}{229}{\boxspacing}
\begin{Verbatim}[commandchars=\\\{\}]
\PY{n}{years}\PY{+w}{ }\PY{o}{=}\PY{+w}{ }\PY{l+m+mi}{1790}\PY{o}{:}\PY{l+m+mi}{10}\PY{o}{:}\PY{l+m+mi}{2010}
\PY{n}{ts}\PY{+w}{ }\PY{o}{=}\PY{+w}{ }\PY{p}{[}\PY{n}{i}\PY{+w}{ }\PY{o}{\PYZhy{}}\PY{+w}{ }\PY{l+m+mi}{1}\PY{+w}{ }\PY{k}{for}\PY{+w}{ }\PY{n}{i}\PY{+w}{ }\PY{k}{in}\PY{+w}{ }\PY{n}{range}\PY{p}{(}\PY{l+m+mi}{1}\PY{p}{,}\PY{+w}{ }\PY{n}{length}\PY{p}{(}\PY{n}{years}\PY{p}{)}\PY{p}{)}\PY{p}{]}\PY{p}{;}
\PY{n}{A}\PY{+w}{ }\PY{o}{=}\PY{+w}{ }\PY{p}{[}\PY{p}{[}\PY{n}{t}\PY{+w}{ }\PY{o}{.\PYZca{}}\PY{+w}{ }\PY{n}{i}\PY{+w}{ }\PY{o}{*}\PY{+w}{ }\PY{l+m+mf}{1.0}\PY{+w}{ }\PY{k}{for}\PY{+w}{ }\PY{n}{t}\PY{+w}{ }\PY{k}{in}\PY{+w}{ }\PY{n}{ts}\PY{p}{]}\PY{+w}{ }\PY{k}{for}\PY{+w}{ }\PY{n}{i}\PY{+w}{ }\PY{k}{in}\PY{+w}{ }\PY{n}{range}\PY{p}{(}\PY{l+m+mi}{0}\PY{p}{,}\PY{l+m+mi}{3}\PY{p}{)}\PY{p}{]}
\PY{n}{A}\PY{+w}{ }\PY{o}{=}\PY{+w}{ }\PY{n}{reshape}\PY{p}{(}\PY{n}{collect}\PY{p}{(}\PY{n}{Iterators}\PY{o}{.}\PY{n}{flatten}\PY{p}{(}\PY{n}{A}\PY{p}{)}\PY{p}{)}\PY{p}{,}\PY{+w}{ }\PY{l+m+mi}{23}\PY{p}{,}\PY{+w}{ }\PY{l+m+mi}{4}\PY{p}{)}

\PY{n}{cs}\PY{+w}{ }\PY{o}{=}\PY{+w}{ }\PY{n}{A}\PY{+w}{ }\PY{o}{\PYZbs{}}\PY{+w}{ }\PY{n}{usa\PYZus{}pop}

\PY{n}{p}\PY{p}{(}\PY{n}{x}\PY{p}{)}\PY{+w}{ }\PY{o}{=}\PY{+w}{ }\PY{n}{sum}\PY{p}{(}\PY{n}{cs}\PY{+w}{ }\PY{o}{.*}\PY{+w}{ }\PY{p}{(}\PY{p}{(}\PY{p}{(}\PY{n}{x}\PY{+w}{ }\PY{o}{\PYZhy{}}\PY{+w}{ }\PY{l+m+mi}{1790}\PY{p}{)}\PY{+w}{ }\PY{o}{/}\PY{+w}{ }\PY{l+m+mi}{10}\PY{p}{)}\PY{+w}{ }\PY{o}{.\PYZca{}}\PY{+w}{ }\PY{p}{[}\PY{l+m+mi}{0}\PY{p}{,}\PY{+w}{ }\PY{l+m+mi}{1}\PY{p}{,}\PY{+w}{ }\PY{l+m+mi}{2}\PY{p}{,}\PY{+w}{ }\PY{l+m+mi}{3}\PY{p}{]}\PY{p}{)}\PY{p}{)}

\PY{n}{println}\PY{p}{(}\PY{l+s}{\PYZdq{}}\PY{l+s}{Predition: }\PY{l+s}{\PYZdq{}}\PY{p}{)}
\PY{n}{println}\PY{p}{(}\PY{l+s}{\PYZdq{}}\PY{l+s}{   2000: }\PY{l+s}{\PYZdq{}}\PY{p}{,}\PY{+w}{ }\PY{n}{p}\PY{p}{(}\PY{l+m+mi}{2000}\PY{p}{)}\PY{p}{,}\PY{+w}{ }\PY{l+s}{\PYZdq{}}\PY{l+s}{ million   Real: }\PY{l+s}{\PYZdq{}}\PY{p}{,}\PY{+w}{ }\PY{n}{usa\PYZus{}pop}\PY{p}{[}\PY{k}{end}\PY{o}{\PYZhy{}}\PY{l+m+mi}{1}\PY{p}{]}\PY{p}{,}\PY{+w}{ }\PY{l+s}{\PYZdq{}}\PY{l+s}{ million}\PY{l+s}{\PYZdq{}}\PY{p}{)}
\PY{n}{println}\PY{p}{(}\PY{l+s}{\PYZdq{}}\PY{l+s}{   2010: }\PY{l+s}{\PYZdq{}}\PY{p}{,}\PY{+w}{ }\PY{n}{p}\PY{p}{(}\PY{l+m+mi}{2010}\PY{p}{)}\PY{p}{,}\PY{+w}{ }\PY{l+s}{\PYZdq{}}\PY{l+s}{ million   Real: }\PY{l+s}{\PYZdq{}}\PY{p}{,}\PY{+w}{ }\PY{n}{usa\PYZus{}pop}\PY{p}{[}\PY{k}{end}\PY{p}{]}\PY{p}{,}\PY{+w}{ }\PY{l+s}{\PYZdq{}}\PY{l+s}{ million}\PY{l+s}{\PYZdq{}}\PY{p}{)}
\PY{n}{println}\PY{p}{(}\PY{l+s}{\PYZdq{}}\PY{l+s}{   2020: }\PY{l+s}{\PYZdq{}}\PY{p}{,}\PY{+w}{ }\PY{n}{p}\PY{p}{(}\PY{l+m+mi}{2020}\PY{p}{)}\PY{p}{,}\PY{+w}{ }\PY{l+s}{\PYZdq{}}\PY{l+s}{ million   Real (From statistics bureau): 3.295 million}\PY{l+s}{\PYZdq{}}\PY{p}{)}

\PY{n}{plot}\PY{p}{(}\PY{n}{p}\PY{p}{,}\PY{+w}{ }\PY{l+m+mi}{1785}\PY{p}{,}\PY{+w}{ }\PY{l+m+mi}{2025}\PY{p}{,}\PY{+w}{ }\PY{n}{xlabel}\PY{o}{=}\PY{l+s}{\PYZdq{}}\PY{l+s}{Year}\PY{l+s}{\PYZdq{}}\PY{p}{,}\PY{+w}{ }\PY{n}{ylabel}\PY{o}{=}\PY{l+s}{\PYZdq{}}\PY{l+s}{Population (million)}\PY{l+s}{\PYZdq{}}\PY{p}{,}
\PY{+w}{    }\PY{n}{title}\PY{o}{=}\PY{l+s}{\PYZdq{}}\PY{l+s}{Population of US versus time}\PY{l+s}{\PYZdq{}}\PY{p}{,}
\PY{+w}{    }\PY{n}{label}\PY{o}{=}\PY{l+s}{\PYZdq{}}\PY{l+s}{Fitting Curve}\PY{l+s}{\PYZdq{}}
\PY{+w}{    }\PY{p}{)}
\PY{n}{scatter!}\PY{p}{(}\PY{n}{years}\PY{p}{,}\PY{+w}{ }\PY{n}{usa\PYZus{}pop}\PY{p}{,}\PY{+w}{ }\PY{n}{label}\PY{o}{=}\PY{l+s}{\PYZdq{}}\PY{l+s}{Samples}\PY{l+s}{\PYZdq{}}\PY{p}{)}
\end{Verbatim}
\end{tcolorbox}

    \begin{Verbatim}[commandchars=\\\{\}]
Predition:
   2000: 279.1552670112495 million   Real: 281.4 million
   2010: 308.53434046822736 million   Real: 308.7 million
   2020: 339.6777617165442 million   Real (From statistics bureau): 3.295
million
    \end{Verbatim}

    \begin{center}
    \adjustimage{max size={0.9\linewidth}{0.9\paperheight}}{hw-02_files/hw-02_49_1.pdf}
    \end{center}
    { \hspace*{\fill} \\}
    
    \begin{enumerate}
\def\labelenumi{(\alph{enumi})}
\setcounter{enumi}{1}
\tightlist
\item
  Look up the actual U.S. population in 2000, 2010, and 2020 and compare
  to the predictions of part (a).
\end{enumerate}

    \begin{longtable}[]{@{}ccc@{}}
\toprule\noalign{}
Year & Prediction (million) & Real (million) \\
\midrule\noalign{}
\endhead
\bottomrule\noalign{}
\endlastfoot
2000 & 279.1552670112495 & 281.4 \\
2010 & 308.53434046822736 & 308.7 \\
2020 & 339.6777617165442 & 329.5 \\
\end{longtable}

From the table, we can note that the predition for years 2000 and 2010
are sufficiently close to the real number (relative error \textless{}
1\%\textgreater), this partly due to that we included these 2 points
when calculating the fitting curve.

For year 2022, the predition is still precise (relative \textless{}
3\%), but the error goes larger.

So we can roughly conclude that the population increses cubicly with
respect to the time.

    \hypertarget{section-3.1-6}{%
\subsection{Section 3.1 \#6}\label{section-3.1-6}}

One series for finding \(\pi\) is

\[\frac{\pi}{2} = 1 + \frac{1}{3} + \frac{1\cdot 2}{3\cdot 5} + \frac{1\cdot 2\cdot 3}{3\cdot 5\cdot 7} + \dots\]

Define \(s_k\) to be the sum of the first \(k\) terms on the right-hand
side, and let \(e_k = |s_k - \pi/2|\).

\begin{enumerate}
\def\labelenumi{(\alph{enumi})}
\tightlist
\item
  Calculate \(e_k\) for \(k=1, \dots, 20\), and plot the sequence on a
  log-linear scale.
\end{enumerate}

    \begin{tcolorbox}[breakable, size=fbox, boxrule=1pt, pad at break*=1mm,colback=cellbackground, colframe=cellborder]
\prompt{In}{incolor}{230}{\boxspacing}
\begin{Verbatim}[commandchars=\\\{\}]
\PY{n}{ks}\PY{+w}{ }\PY{o}{=}\PY{+w}{ }\PY{l+m+mi}{1}\PY{o}{:}\PY{l+m+mi}{20}

\PY{n}{cur}\PY{+w}{ }\PY{o}{=}\PY{+w}{ }\PY{l+m+mi}{1}
\PY{n}{ss}\PY{+w}{ }\PY{o}{=}\PY{+w}{ }\PY{p}{[}\PY{n}{cur}\PY{p}{]}
\PY{k}{for}\PY{+w}{ }\PY{n}{k}\PY{+w}{ }\PY{k}{in}\PY{+w}{ }\PY{n}{ks}\PY{p}{[}\PY{l+m+mi}{2}\PY{o}{:}\PY{k}{end}\PY{p}{]}\PY{p}{;}
\PY{+w}{    }\PY{n}{cur}\PY{+w}{ }\PY{o}{*=}\PY{+w}{ }\PY{p}{(}\PY{n}{k}\PY{+w}{ }\PY{o}{\PYZhy{}}\PY{+w}{ }\PY{l+m+mi}{1}\PY{p}{)}\PY{+w}{ }\PY{o}{/}\PY{+w}{ }\PY{p}{(}\PY{l+m+mi}{2}\PY{+w}{ }\PY{o}{*}\PY{+w}{ }\PY{n}{k}\PY{+w}{ }\PY{o}{\PYZhy{}}\PY{+w}{ }\PY{l+m+mi}{1}\PY{p}{)}
\PY{+w}{    }\PY{n}{ss}\PY{+w}{ }\PY{o}{=}\PY{+w}{ }\PY{p}{[}\PY{n}{ss}\PY{+w}{ }\PY{n}{ss}\PY{p}{[}\PY{k}{end}\PY{p}{]}\PY{o}{+}\PY{n}{cur}\PY{p}{]}
\PY{k}{end}
\PY{n}{ss}\PY{+w}{ }\PY{o}{=}\PY{+w}{ }\PY{n}{ss}\PY{p}{[}\PY{o}{:}\PY{p}{]}

\PY{n}{es}\PY{+w}{ }\PY{o}{=}\PY{+w}{ }\PY{p}{[}\PY{n}{abs}\PY{p}{(}\PY{n}{s}\PY{+w}{ }\PY{o}{\PYZhy{}}\PY{+w}{ }\PY{n+nb}{pi}\PY{+w}{ }\PY{o}{/}\PY{+w}{ }\PY{l+m+mi}{2}\PY{p}{)}\PY{+w}{ }\PY{k}{for}\PY{+w}{ }\PY{n}{s}\PY{+w}{ }\PY{k}{in}\PY{+w}{ }\PY{n}{ss}\PY{p}{]}
\end{Verbatim}
\end{tcolorbox}

    
    \begin{Verbatim}[commandchars=\\\{\}]
20-element Vector\{Float64\}:
 0.5707963267948966
 0.2374629934615633
 0.10412966012823
 0.04698680298537283
 0.021589977588547526
 0.010045966044536003
 0.004717960716530634
 0.0022315582301282166
 0.0010614864718212491
 0.0005072419547285101
 0.00024331599420812644
 0.00011709053482888976
 5.650231432685615e-5
 2.7330208159259683e-5
 1.3247122423232227e-5
 6.432726099347974e-6
 3.1287763664344936e-6
 1.5240007820160884e-6
 7.432991462330563e-7
 3.6295732375180023e-7
    \end{Verbatim}

    
    \begin{tcolorbox}[breakable, size=fbox, boxrule=1pt, pad at break*=1mm,colback=cellbackground, colframe=cellborder]
\prompt{In}{incolor}{231}{\boxspacing}
\begin{Verbatim}[commandchars=\\\{\}]
\PY{n}{plot}\PY{p}{(}\PY{n}{es}\PY{p}{,}\PY{+w}{ }\PY{n}{yaxis}\PY{o}{=}\PY{l+s+ss}{:log}\PY{p}{,}\PY{+w}{ }\PY{n}{label}\PY{o}{=}\PY{l+s}{\PYZdq{}}\PY{l+s}{Absolute Error}\PY{l+s}{\PYZdq{}}\PY{p}{,}\PY{+w}{ }\PY{n}{ylabel}\PY{o}{=}\PY{l+s}{\PYZdq{}}\PY{l+s}{Error}\PY{l+s}{\PYZdq{}}\PY{p}{,}\PY{+w}{ }\PY{n}{xlabel}\PY{o}{=}\PY{l+s}{\PYZdq{}}\PY{l+s}{\PYZsh{}Items}\PY{l+s}{\PYZdq{}}\PY{p}{,}
\PY{+w}{    }\PY{n}{title}\PY{o}{=}\PY{l+s}{\PYZdq{}}\PY{l+s}{Absolute errors of approxmation |s\PYZus{}k \PYZhy{} pi/2|}\PY{l+s}{\PYZdq{}}\PY{p}{)}
\end{Verbatim}
\end{tcolorbox}

    \begin{center}
    \adjustimage{max size={0.9\linewidth}{0.9\paperheight}}{hw-02_files/hw-02_54_0.pdf}
    \end{center}
    { \hspace*{\fill} \\}
    
    \begin{enumerate}
\def\labelenumi{(\alph{enumi})}
\setcounter{enumi}{1}
\tightlist
\item
  Determine \(a\) and \(b\) in a least-squares fit
  \(e_k \approx a\cdot b^k\), and superimpose the fit on the plot from
  part (a).
\end{enumerate}

Hints:

\begin{itemize}
\tightlist
\item
  It's easy to make unnecessary errors in the terms of this series as
  \(k\) gets large (about 20). Be careful how you calculate the terms in
  the sum.
\item
  Use \texttt{yaxis\ =\ :log} to transform the axis and keep pretty
  labels.
\end{itemize}

    \begin{tcolorbox}[breakable, size=fbox, boxrule=1pt, pad at break*=1mm,colback=cellbackground, colframe=cellborder]
\prompt{In}{incolor}{232}{\boxspacing}
\begin{Verbatim}[commandchars=\\\{\}]
\PY{k}{using}\PY{+w}{ }\PY{n}{Plots}

\PY{n}{A}\PY{+w}{ }\PY{o}{=}\PY{+w}{ }\PY{p}{[}\PY{n}{ones}\PY{p}{(}\PY{n}{length}\PY{p}{(}\PY{n}{ks}\PY{p}{)}\PY{p}{)}\PY{+w}{ }\PY{n}{ks}\PY{p}{]}

\PY{n}{cs}\PY{+w}{ }\PY{o}{=}\PY{+w}{ }\PY{n}{A}\PY{+w}{ }\PY{o}{\PYZbs{}}\PY{+w}{ }\PY{n}{log}\PY{o}{.}\PY{p}{(}\PY{n}{es}\PY{p}{)}
\PY{n}{a}\PY{+w}{ }\PY{o}{=}\PY{+w}{ }\PY{n}{exp}\PY{p}{(}\PY{n}{cs}\PY{p}{[}\PY{l+m+mi}{1}\PY{p}{]}\PY{p}{)}\PY{p}{;}\PY{+w}{ }\PY{n}{b}\PY{+w}{ }\PY{o}{=}\PY{+w}{ }\PY{n}{exp}\PY{p}{(}\PY{n}{cs}\PY{p}{[}\PY{l+m+mi}{2}\PY{p}{]}\PY{p}{)}

\PY{n}{ys}\PY{+w}{ }\PY{o}{=}\PY{+w}{ }\PY{p}{[}\PY{n}{a}\PY{+w}{ }\PY{o}{*}\PY{+w}{ }\PY{n}{b}\PY{+w}{ }\PY{o}{\PYZca{}}\PY{+w}{ }\PY{n}{k}\PY{+w}{ }\PY{k}{for}\PY{+w}{ }\PY{n}{k}\PY{+w}{ }\PY{k}{in}\PY{+w}{ }\PY{n}{ks}\PY{p}{]}

\PY{n}{ekf}\PY{p}{(}\PY{n}{k}\PY{p}{)}\PY{+w}{ }\PY{o}{=}\PY{+w}{ }\PY{n}{a}\PY{+w}{ }\PY{o}{*}\PY{+w}{ }\PY{n}{b}\PY{+w}{ }\PY{o}{\PYZca{}}\PY{+w}{ }\PY{n}{k}

\PY{n}{println}\PY{p}{(}\PY{l+s}{\PYZdq{}}\PY{l+s}{a: }\PY{l+s}{\PYZdq{}}\PY{p}{,}\PY{+w}{ }\PY{n}{a}\PY{p}{,}\PY{+w}{ }\PY{l+s}{\PYZdq{}}\PY{l+s}{    b: }\PY{l+s}{\PYZdq{}}\PY{p}{,}\PY{+w}{ }\PY{n}{b}\PY{p}{)}

\PY{n}{plot}\PY{p}{(}\PY{n}{ekf}\PY{p}{,}\PY{+w}{ }\PY{l+m+mi}{0}\PY{p}{,}\PY{+w}{ }\PY{l+m+mi}{21}\PY{p}{,}\PY{+w}{ }\PY{n}{yaxis}\PY{o}{=}\PY{l+s+ss}{:log}\PY{p}{,}\PY{+w}{ }
\PY{+w}{    }\PY{n}{xlabel}\PY{o}{=}\PY{l+s}{\PYZdq{}}\PY{l+s}{\PYZsh{}Items}\PY{l+s}{\PYZdq{}}\PY{p}{,}\PY{+w}{ }\PY{n}{ylabel}\PY{o}{=}\PY{l+s}{\PYZdq{}}\PY{l+s}{Absolute Errors}\PY{l+s}{\PYZdq{}}\PY{p}{,}
\PY{+w}{    }\PY{n}{label}\PY{o}{=}\PY{l+s}{\PYZdq{}}\PY{l+s}{Fitting Curve}\PY{l+s}{\PYZdq{}}\PY{p}{,}\PY{+w}{ }\PY{n}{title}\PY{o}{=}\PY{l+s}{\PYZdq{}}\PY{l+s}{Absolution errors versus number of items}\PY{l+s}{\PYZdq{}}\PY{p}{)}
\PY{n}{scatter!}\PY{p}{(}\PY{n}{ks}\PY{p}{,}\PY{+w}{ }\PY{n}{ys}\PY{p}{,}\PY{+w}{ }\PY{n}{label}\PY{o}{=}\PY{l+s}{\PYZdq{}}\PY{l+s}{Superimposed Points}\PY{l+s}{\PYZdq{}}\PY{p}{)}
\PY{n}{scatter!}\PY{p}{(}\PY{n}{ks}\PY{p}{,}\PY{+w}{ }\PY{n}{es}\PY{p}{,}\PY{+w}{ }\PY{n}{label}\PY{o}{=}\PY{l+s}{\PYZdq{}}\PY{l+s}{Data Points}\PY{l+s}{\PYZdq{}}\PY{p}{)}
\end{Verbatim}
\end{tcolorbox}

    \begin{Verbatim}[commandchars=\\\{\}]
a: 0.9371482082402564    b: 0.47527526723108526
    \end{Verbatim}

    \begin{center}
    \adjustimage{max size={0.9\linewidth}{0.9\paperheight}}{hw-02_files/hw-02_56_1.pdf}
    \end{center}
    { \hspace*{\fill} \\}
    

    % Add a bibliography block to the postdoc
    
    
    
\end{document}
